\begin{register}{H}{PMU\_HP\_ACTIVE\_DIG\_POWER\_REG}{0x0000}\label{regdesc:PMUHPACTIVEDIGPOWERREG}
\iftagged{ESP32-P4-latest}{
    \regfield{PMU\_HP\_ACTIVE\_PD\_TOP\_PD\_EN}{1}{31}{{0}}%
    \regfield{PMU\_HP\_ACTIVE\_PD\_CNNT\_PD\_EN}{1}{30}{{0}}%
    \regfield{PMU\_HP\_ACTIVE\_PD\_HP\_CPU\_PD\_EN}{1}{29}{{0}}%
    \regfield{(reserved)}{5}{24}{00000}%
    \regfield{PMU\_HP\_ACTIVE\_PD\_HP\_MEM\_PD\_EN}{1}{23}{{0}}%
    \regfield{PMU\_HP\_ACTIVE\_HP\_MEM\_DSLP}{1}{22}{{0}}%
    \regfield{PMU\_HP\_ACTIVE\_DCDC\_SWITCH\_PD\_EN}{1}{21}{{0}}%
    \regfield{(reserved)}{21}{0}{000000000000000000000}%
    \reglabel{Reset}\regnewline%
}{
    \regfield{PMU\_HP\_ACTIVE\_PD\_TOP\_PD\_EN}{1}{31}{{0}}%
    \regfield{PMU\_HP\_ACTIVE\_PD\_CNNT\_PD\_EN}{1}{30}{{0}}%
    \regfield{(reserved)}{6}{24}{000000}%
    \regfield{PMU\_HP\_ACTIVE\_PD\_HP\_MEM\_PD\_EN}{1}{23}{{0}}%
    \regfield{PMU\_HP\_ACTIVE\_HP\_MEM\_DSLP}{1}{22}{{0}}%
    \regfield{PMU\_HP\_ACTIVE\_DCDC\_SWITCH\_PD\_EN}{1}{21}{{0}}%
    \regfield{(reserved)}{21}{0}{000000000000000000000}%
    \reglabel{Reset}\regnewline%
}
\begin{regdesc}\begin{reglist}
\label{fielddesc:PMUHPACTIVEDCDCSWITCHPDEN}\item [PMU\_HP\_ACTIVE\_DCDC\_SWITCH\_PD\_EN] 配置是否在 HP\_ACTIVE 状态下断开所有来源于外部 DCDC 的供电。\\
0：不断开 \\
1：断开 \\
(R/W)
\label{fielddesc:PMUHPACTIVEHPMEMDSLP}\item [PMU\_HP\_ACTIVE\_HP\_MEM\_DSLP] 配置是否在 HP\_ACTIVE 状态下让 L2MEM\_G0–G5 进入 Deep-sleep。\\
0：不进入 Deep-sleep \\
1：进入 Deep-sleep\\
(R/W)
\label{fielddesc:PMUHPACTIVEPDHPMEMPDEN}\item [PMU\_HP\_ACTIVE\_PD\_HP\_MEM\_PD\_EN] 配置是否在 HP\_ACTIVE 状态下关闭 L2MEM\_G0–G5 电源域。 \\
0：打开\\
1：关闭 \\
(R/W)

\tagged{ESP32-P4-latest}{
    \label{fielddesc:PMUHPACTIVEPDHPCPUPDEN}\item [PMU\_HP\_ACTIVE\_PD\_HP\_CPU\_PD\_EN] 配置是否在 HP\_ACTIVE 状态下关闭 HP\_CPU 电源域。 \\
    0：打开\\
    1：关闭 \\
    (R/W)
}

\label{fielddesc:PMUHPACTIVEPDCNNTPDEN}\item [PMU\_HP\_ACTIVE\_PD\_CNNT\_PD\_EN] 配置是否在 HP\_ACTIVE 状态下关闭 HP\_CNNT 电源域。 \\
0：打开\\
1：关闭 \\
(R/W)
\label{fielddesc:PMUHPACTIVEPDTOPPDEN}\item [PMU\_HP\_ACTIVE\_PD\_TOP\_PD\_EN] 配置是否在 HP\_ACTIVE 状态下关闭 PD\_TOP 电源域。\\
0：打开\\
1：关闭 \\
(R/W)
\end{reglist}\end{regdesc}
\end{register}


\begin{register}{H}{PMU\_HP\_ACTIVE\_ICG\_HP\_FUNC\_REG}{0x0004}\label{regdesc:PMUHPACTIVEICGHPFUNCREG}
\regfield{PMU\_HP\_ACTIVE\_DIG\_ICG\_FUNC\_EN}{32}{0}{{0xffffffff}}%
\reglabel{Reset}\regnewline%
\begin{regdesc}\begin{reglist}
\label{fielddesc:PMUHPACTIVEDIGICGFUNCEN}\item [PMU\_HP\_ACTIVE\_DIG\_ICG\_FUNC\_EN] 配置 HP\_ACTIVE 状态下高性能系统外设功能时钟的开关。具体配置请见表 \ref{tab:lowpowm-hp-sys-peri-clk}。 \\
0：关闭\\
1：打开\\
(R/W)
\end{reglist}\end{regdesc}
\end{register}


\begin{register}{H}{PMU\_HP\_ACTIVE\_HP\_SYS\_CNTL\_REG}{0x0010}\label{regdesc:PMUHPACTIVEHPSYSCNTLREG}
\regfield{(reserved)}{2}{30}{00}%
\regfield{PMU\_HP\_ACTIVE\_DIG\_CPU\_STALL}{1}{29}{{0}}%
\regfield{PMU\_HP\_ACTIVE\_DIG\_PAUSE\_WDT}{1}{28}{{0}}%
\regfield{PMU\_HP\_ACTIVE\_DIG\_PAD\_SLP\_SEL}{1}{27}{{0}}%
\regfield{PMU\_HP\_ACTIVE\_HP\_PAD\_HOLD\_ALL}{1}{26}{{0}}%
\regfield{PMU\_HP\_ACTIVE\_LP\_PAD\_HOLD\_ALL}{1}{25}{{0}}%
\regfield{PMU\_HP\_ACTIVE\_UART\_WAKEUP\_EN}{1}{24}{{0}}%
\regfield{(reserved)}{24}{0}{000000000000000000000000}%
\reglabel{Reset}\regnewline%
\begin{regdesc}\begin{reglist}
\label{fielddesc:PMUHPACTIVEUARTWAKEUPEN}\item [PMU\_HP\_ACTIVE\_UART\_WAKEUP\_EN] 配置是否在 HP\_ACTIVE 状态下使能 UART 唤醒功能。 \\
0：关闭唤醒功能\\
1：使能唤醒功能\\
(R/W)
\label{fielddesc:PMUHPACTIVELPPADHOLDALL}\item [PMU\_HP\_ACTIVE\_LP\_PAD\_HOLD\_ALL] 配置是否在 HP\_ACTIVE 状态下保持 LP GPIO 的配置。\\
0：不保持\\
1：保持\\
(R/W)
\label{fielddesc:PMUHPACTIVEHPPADHOLDALL}\item [PMU\_HP\_ACTIVE\_HP\_PAD\_HOLD\_ALL] 配置是否在 HP\_ACTIVE 状态下保持 GPIO 的配置。\\
0：不保持\\
1：保持\\
(R/W)
\label{fielddesc:PMUHPACTIVEDIGPADSLPSEL}\item [PMU\_HP\_ACTIVE\_DIG\_PAD\_SLP\_SEL] 配置是否在 HP\_ACTIVE 状态下使用 GPIO 的 Light-sleep 模式配置。\\
0：正常工作模式的配置\\
1：Light-sleep 模式配置，详见章节 \ref{mod:iomuxgpio} \textit{\nameref{mod:iomuxgpio}} > 小节 \ref{sec:iomuxgpio-pin-func-light-sleep} \textit{\nameref{sec:iomuxgpio-pin-func-light-sleep}}。\\
(R/W)
\label{fielddesc:PMUHPACTIVEDIGPAUSEWDT}\item [PMU\_HP\_ACTIVE\_DIG\_PAUSE\_WDT] 配置是否在 HP\_ACTIVE 状态下暂停看门狗。 \\
0：不暂停\\
1：暂停\\
(R/W)
\label{fielddesc:PMUHPACTIVEDIGCPUSTALL}\item [PMU\_HP\_ACTIVE\_DIG\_CPU\_STALL] 配置是否在 HP\_ACTIVE 状态下暂停 CPU。 \\
0：不暂停\\
1：暂停\\
(R/W)
\end{reglist}\end{regdesc}
\end{register}


\begin{register}{H}{PMU\_HP\_ACTIVE\_HP\_CK\_POWER\_REG}{0x0014}\label{regdesc:PMUHPACTIVEHPCKPOWERREG}
\regfield{(reserved)}{1}{31}{0}%
\regfield{PMU\_HP\_ACTIVE\_XPD\_PLL}{4}{27}{{0}}%
\regfield{PMU\_HP\_ACTIVE\_XPD\_PLL\_I2C}{4}{23}{{0}}%
\regfield{PMU\_HP\_ACTIVE\_I2C\_RETENTION}{1}{22}{{0}}%
\regfield{PMU\_HP\_ACTIVE\_I2C\_ISO\_EN}{1}{21}{{0}}%
\regfield{(reserved)}{21}{0}{000000000000000000000}%
\reglabel{Reset}\regnewline%
\begin{regdesc}\begin{reglist}
\label{fielddesc:PMUHPACTIVEI2CISOEN}\item [PMU\_HP\_ACTIVE\_I2C\_ISO\_EN]配置是否在 HP\_ACTIVE 状态打开 ANALOG\_I2C。 \\
0：关闭\\
1：开启\\
(R/W)
\label{fielddesc:PMUHPACTIVEI2CRETENTION}\item [PMU\_HP\_ACTIVE\_I2C\_RETENTION] 配置是否在 HP\_ACTIVE 状态使能 ANALOG\_I2C 进入 retention 状态。\\
0：不使能\\
1：使能\\
(R/W)
\label{fielddesc:PMUHPACTIVEXPDPLLI2C}\item [PMU\_HP\_ACTIVE\_XPD\_PLL\_I2C] 配置是否在 HP\_ACTIVE 状态下开启 PLL 的 I2C 控制器，该组寄存器为 bitmap，每个 bit 控制一个 PLL 的 I2C 寄存器组。\\
0：关闭\\
1：开启\\
(R/W)
\label{fielddesc:PMUHPACTIVEXPDPLL}\item [PMU\_HP\_ACTIVE\_XPD\_PLL] 配置是否在 HP\_ACTIVE 状态下开启 HP 系统的模拟电源域，每个bit 控制一个电源域，bit 0 对应 CPLL\_CLK，bit 1 对应 SPLL\_CLK，bit 2 对应 Audio PLL，bit 3 对应 SDIO PLL。\\
0：关闭\\
1：开启\\
(R/W)
\end{reglist}\end{regdesc}
\end{register}


\begin{register}{H}{PMU\_HP\_ACTIVE\_BIAS\_REG}{0x0018}\label{regdesc:PMUHPACTIVEBIASREG}
\regfield{(reserved)}{7}{25}{0000000}%
\regfield{PMU\_HP\_ACTIVE\_DCM\_MODE}{2}{23}{{0}}%
\regfield{PMU\_HP\_ACTIVE\_DCM\_VSET}{5}{18}{{20}}%
\regfield{(reserved)}{18}{0}{000000000000000000}%
\reglabel{Reset}\regnewline%
\begin{regdesc}\begin{reglist}
\label{fielddesc:PMUHPACTIVEDCMVSET}\item [PMU\_HP\_ACTIVE\_DCM\_VSET]  在 HP\_ACTIVE 状态下调节 DCDC 的电压。(R/W)   \label{fielddesc:PMUHPACTIVEDCMMODE}\item [PMU\_HP\_ACTIVE\_DCM\_MODE] 配置是否在 HP\_ACTIVE 状态下使能 DCDC 的状态。\\
0：不使能 \\
1：使能 \\
2, 3：保留\\
(R/W)
\end{reglist}\end{regdesc}
\end{register}


\begin{register}{H}{PMU\_HP\_ACTIVE\_BACKUP\_REG}{0x001C}\label{regdesc:PMUHPACTIVEBACKUPREG}
\regfield{(reserved)}{2}{30}{00}%
\regfield{PMU\_HP\_SLEEP2ACTIVE\_BACKUP\_EN}{1}{29}{{0}}%
\regfield{(reserved)}{3}{26}{000}%
\regfield{(reserved)}{3}{23}{000}%
\regfield{PMU\_HP\_SLEEP2ACTIVE\_BACKUP\_MODE}{3}{20}{{0}}%
\regfield{(reserved)}{2}{18}{00}%
\regfield{(reserved)}{2}{16}{00}%
\regfield{PMU\_HP\_SLEEP2ACTIVE\_BACKUP\_CLK\_SEL}{2}{14}{{0}}%
\regfield{(reserved)}{1}{13}{0}%
\regfield{(reserved)}{1}{12}{0}%
\regfield{(reserved)}{1}{11}{0}%
\regfield{(reserved)}{2}{8}{00}%
\regfield{(reserved)}{2}{6}{{0}}%
\regfield{(reserved)}{2}{4}{{0}}%
\regfield{(reserved)}{4}{0}{0000}%
\reglabel{Reset}\regnewline%
\begin{regdesc}\begin{reglist}
\label{fielddesc:PMUHPSLEEP2ACTIVEBACKUPCLKSEL}\item [PMU\_HP\_SLEEP2ACTIVE\_BACKUP\_CLK\_SEL] 配置备份模块从 HP\_SLEEP 切换为 HP\_ACTIVE 时的功能时钟源。\\
0：时钟源为 XTAL\\
1：时钟源为 PLL\_CLK\\
2：时钟源为 RC\_FAST\_CLK\\
3：无效配置\\
(R/W)
\label{fielddesc:PMUHPSLEEP2ACTIVEBACKUPMODE}\item [PMU\_HP\_SLEEP2ACTIVE\_BACKUP\_MODE] 配置从 HP\_SLEEP 切换为 HP\_ACTIVE 时数据搬移的方向和链表指针。\\
最高位：\\
0：存储器到外设\\
1：外设到存储器\\
低两位：\\
0：PAU\_LINK\_ADDR\_0 \\
1：PAU\_LINK\_ADDR\_1 \\
2：PAU\_LINK\_ADDR\_2 \\
3：PAU\_LINK\_ADDR\_3 \\
(R/W)

\label{fielddesc:PMUHPSLEEP2ACTIVEBACKUPEN}\item [PMU\_HP\_SLEEP2ACTIVE\_BACKUP\_EN] 配置是否在从 HP\_SLEEP 切换为 HP\_ACTIVE 时使能数据搬移。\\
0：禁能\\
1：使能\\
(R/W)
\end{reglist}\end{regdesc}
\end{register}


\begin{register}{H}{PMU\_HP\_ACTIVE\_BACKUP\_CLK\_REG}{0x0020}\label{regdesc:PMUHPACTIVEBACKUPCLKREG}
\regfield{PMU\_HP\_ACTIVE\_BACKUP\_ICG\_FUNC\_EN}{32}{0}{{0}}%
\reglabel{Reset}\regnewline%
\begin{regdesc}\begin{reglist}
\label{fielddesc:PMUHPACTIVEBACKUPICGFUNCEN}\item [PMU\_HP\_ACTIVE\_BACKUP\_ICG\_FUNC\_EN] 配置目标状态为 HP\_ACTIVE 时不同外设功能时钟的开关情况，不同位代表不同外设，详见表 \ref{tab:lowpowm-hp-sys-peri-clk}。\\
0：关闭\\
1：打开\\
(R/W)
\end{reglist}\end{regdesc}
\end{register}


\begin{register}{H}{PMU\_HP\_ACTIVE\_SYSCLK\_REG}{0x0024}\label{regdesc:PMUHPACTIVESYSCLKREG}
\regfield{PMU\_HP\_ACTIVE\_DIG\_SYS\_CLK\_SEL}{2}{30}{{0}}%
\regfield{PMU\_HP\_ACTIVE\_ICG\_SLP\_SEL}{1}{29}{{0}}%
\regfield{PMU\_HP\_ACTIVE\_SYS\_CLK\_SLP\_SEL}{1}{28}{{0}}%
\regfield{PMU\_HP\_ACTIVE\_ICG\_SYS\_CLOCK\_EN}{1}{27}{{0}}%
\regfield{(reserved)}{27}{0}{00000000000000000000000000}%
\reglabel{Reset}\regnewline%
\begin{regdesc}\begin{reglist}
\label{fielddesc:PMUHPACTIVEICGSYSCLOCKEN}\item [PMU\_HP\_ACTIVE\_ICG\_SYS\_CLOCK\_EN] 配置是否在 HP\_ACTIVE 状态下打开 ROOT\_CLK。\\
0：关闭\\
1：打开\\
(R/W)
\label{fielddesc:PMUHPACTIVESYSCLKSLPSEL}\item [PMU\_HP\_ACTIVE\_SYS\_CLK\_SLP\_SEL] 配置是否在 HP\_ACTIVE 状态下由 PMU 控制时钟源。\\
0：由 PCR 寄存器控制\\
1：由 PMU 控制\\
(R/W)
\label{fielddesc:PMUHPACTIVEICGSLPSEL}\item [PMU\_HP\_ACTIVE\_ICG\_SLP\_SEL] 配置是否在 HP\_ACTIVE 状态下是否由 PMU 控制各时钟门控的开关。\\
0：由 PCR 寄存器控制\\
1：由 PMU 控制\\
(R/W)
\label{fielddesc:PMUHPACTIVEDIGSYSCLKSEL}\item [PMU\_HP\_ACTIVE\_DIG\_SYS\_CLK\_SEL] 配置 HP\_ACTIVE 状态下 ROOT\_CLK 的时钟源。\\
0：时钟源为 XTAL\\
1：时钟源为 PLL\_CLK\\
2：时钟源为 RC\_FAST\_CLK\\
3：无效配置\\
(R/W)
\end{reglist}\end{regdesc}
\end{register}


\begin{register}{H}{PMU\_HP\_ACTIVE\_HP\_REGULATOR0\_REG}{0x0028}\label{regdesc:PMUHPACTIVEHPREGULATOR0REG}    \regfield{PMU\_HP\_ACTIVE\_HP\_REGULATOR\_DBIAS}{5}{27}{{24}}
\regfield{(reserved)}{4}{23}{0000}
 \regfield{PMU\_HP\_ACTIVE\_REGULATOR\_VO1\_XPD}{1}{22}{{1}}%
\regfield{(reserved)}{3}{19}{000}
\regfield{PMU\_HP\_ACTIVE\_HP\_REGULATOR\_XPD}{1}{18}{{1}}%
\regfield{(reserved)}{2}{16}{00}%
\regfield{PMU\_DIG\_DBIAS\_INIT}{1}{15}{{0}}
\regfield{PMU\_DIG\_REGULATOR0\_DBIAS\_SEL}{1}{14}{{1}}%
\regfield{PMU\_HP\_DBIAS\_VOL}{5}{9}{{24}}%
\regfield{PMU\_LP\_DBIAS\_VOL}{5}{4}{{24}}%
\regfield{(reserved)}{4}{0}{0000}%
\reglabel{Reset}\regnewline%
\begin{regdesc}\begin{reglist}
\label{fielddesc:PMULPDBIASVOL}\item [PMU\_LP\_DBIAS\_VOL] 表示当前的低功耗系统调压器的电压。(RO)
\label{fielddesc:PMUHPDBIASVOL}\item [PMU\_HP\_DBIAS\_VOL] 表示当前的高性能系统调压器的电压。(RO)
\label{fielddesc:PMUDIGREGULATOR0DBIASSEL}\item [PMU\_DIG\_REGULATOR0\_DBIAS\_SEL] 选择高性能系统调压器的电压控制权。\\
0：硬件自动调压 \\
1：软件调压 \\
(R/W)    \label{fielddesc:PMUDIGDBIASINIT}\item [PMU\_DIG\_DBIAS\_INIT] 更新初始化 PVT 的电压配置。 (WT)
\label{fielddesc:PMUHPACTIVEHPREGULATORXPD}\item [PMU\_HP\_ACTIVE\_HP\_REGULATOR\_XPD] 配置是否在 HP\_ACTIVE 状态下打开高性能系统调压器。\\
0：关闭高性能系统调压器 \\
1：打开高性能系统调压器\\
(R/W)
\label{fielddesc:PMUHPACTIVEREGULATORVO1XPD}\item[PMU\_HP\_ACTIVE\_REGULATOR\_VO1\_XPD] 配置 HP\_ACTIVE 状态下 VO1 的开关。 \\
0：关闭 \\
1：打开 \\
(R/W)
\label{fielddesc:PMUHPACTIVEHPREGULATORDBIAS}\item [PMU\_HP\_ACTIVE\_HP\_REGULATOR\_DBIAS] 在 HP\_ACTIVE 状态下调节高性能系统调压器的电压，寄存器值越大，电压越高。(R/W)
\end{reglist}\end{regdesc}
\end{register}



\begin{register}{H}{PMU\_HP\_ACTIVE\_XTAL\_REG}{0x0030}\label{regdesc:PMUHPACTIVEXTALREG}
\regfield{PMU\_HP\_ACTIVE\_XPD\_XTAL}{1}{31}{{1}}%
\regfield{(reserved)}{31}{0}{0000000000000000000000000000000}%
\reglabel{Reset}\regnewline%
\begin{regdesc}\begin{reglist}
\label{fielddesc:PMUHPACTIVEXPDXTAL}\item [PMU\_HP\_ACTIVE\_XPD\_XTAL] 配置是否在 HP\_ACTIVE 状态下打开 XTAL\_CLK 的模拟源。\\
0：关闭\\
1：打开\\
(R/W)
\end{reglist}\end{regdesc}
\end{register}


\begin{register}{H}{PMU\_HP\_SLEEP\_DIG\_POWER\_REG}{0x0068}\label{regdesc:PMUHPSLEEPDIGPOWERREG}
\iftagged{ESP32-P4-latest}{
    \regfield{PMU\_HP\_SLEEP\_PD\_TOP\_PD\_EN}{1}{31}{{0}}%
    \regfield{PMU\_HP\_SLEEP\_PD\_CNNT\_PD\_EN}{1}{30}{{0}}%
    \regfield{PMU\_HP\_SLEEP\_PD\_HP\_CPU\_PD\_EN}{1}{29}{{0}}%
    \regfield{(reserved)}{5}{24}{00000}%
    \regfield{PMU\_HP\_SLEEP\_PD\_HP\_MEM\_PD\_EN}{1}{23}{{0}}%
    \regfield{PMU\_HP\_SLEEP\_HP\_MEM\_DSLP}{1}{22}{{0}}%
    \regfield{PMU\_HP\_SLEEP\_DCDC\_SWITCH\_PD\_EN}{1}{21}{{0}}%
    \regfield{(reserved)}{21}{0}{000000000000000000000}%
    \reglabel{Reset}\regnewline%
}{
    \regfield{PMU\_HP\_SLEEP\_PD\_TOP\_PD\_EN}{1}{31}{{0}}%
    \regfield{PMU\_HP\_SLEEP\_PD\_CNNT\_PD\_EN}{1}{30}{{0}}%
    \regfield{(reserved)}{6}{24}{000000}%
    \regfield{PMU\_HP\_SLEEP\_PD\_HP\_MEM\_PD\_EN}{1}{23}{{0}}%
    \regfield{PMU\_HP\_SLEEP\_HP\_MEM\_DSLP}{1}{22}{{0}}%
    \regfield{PMU\_HP\_SLEEP\_DCDC\_SWITCH\_PD\_EN}{1}{21}{{0}}%
    \regfield{(reserved)}{21}{0}{000000000000000000000}%
    \reglabel{Reset}\regnewline%
}
\begin{regdesc}\begin{reglist}
\label{fielddesc:PMUHPSLEEPDCDCSWITCHPDEN}\item [PMU\_HP\_SLEEP\_DCDC\_SWITCH\_PD\_EN] 配置是否在 HP\_SLEEP 状态下断开所有来源于外部 DCDC 的供电。\\
0：不断开 \\
1：断开 \\
(R/W)
\label{fielddesc:PMUHPSLEEPHPMEMDSLP}\item [PMU\_HP\_SLEEP\_HP\_MEM\_DSLP] 配置是否在 HP\_SLEEP 状态下让 L2MEM\_G0–G5 进入 Deep-sleep。\\
0：不进入 Deep-sleep \\
1：进入 Deep-sleep\\
(R/W)
\label{fielddesc:PMUHPSLEEPPDHPMEMPDEN}\item [PMU\_HP\_SLEEP\_PD\_HP\_MEM\_PD\_EN] 配置是否在 HP\_SLEEP 状态下关闭 L2MEM\_G0–G5 电源域。 \\
0：打开\\
1：关闭 \\
(R/W)

\tagged{ESP32-P4-latest}{
    \label{fielddesc:PMUHPSLEEPPDHPCPUPDEN}\item [PMU\_HP\_SLEEP\_PD\_HP\_CPU\_PD\_EN] 配置是否在 HP\_SLEEP 状态下关闭 HP\_CPU 电源域 \\
    0：打开\\
    1：关闭 \\
    (R/W)
}
\label{fielddesc:PMUHPSLEEPPDCNNTPDEN}\item [PMU\_HP\_SLEEP\_PD\_CNNT\_PD\_EN] 配置是否在 HP\_SLEEP 状态下关闭 HP\_CNNT 电源域 \\
0：打开\\
1：关闭 \\
(R/W)
\label{fielddesc:PMUHPSLEEPPDTOPPDEN}\item [PMU\_HP\_SLEEP\_PD\_TOP\_PD\_EN] 配置是否在 HP\_SLEEP 状态下关闭 PD\_TOP 电源域。\\
0：打开\\
1：关闭 \\
(R/W)
\end{reglist}\end{regdesc}
\end{register}

\begin{register}{H}{PMU\_HP\_SLEEP\_ICG\_HP\_FUNC\_REG}{0x006C}\label{regdesc:PMUHPSLEEPICGHPFUNCREG}
\regfield{PMU\_HP\_SLEEP\_DIG\_ICG\_FUNC\_EN}{32}{0}{{0xffffffff}}%
\reglabel{Reset}\regnewline%
\begin{regdesc}\begin{reglist}
\label{fielddesc:PMUHPSLEEPDIGICGFUNCEN}\item [PMU\_HP\_SLEEP\_DIG\_ICG\_FUNC\_EN] 配置 HP\_SLEEP 状态下高性能系统外设功能时钟的开关。具体配置请见表 \ref{tab:lowpowm-hp-sys-peri-clk}。 \\
0：关闭\\
1：打开\\
(R/W)
\end{reglist}\end{regdesc}
\end{register}


\begin{register}{H}{PMU\_HP\_SLEEP\_HP\_SYS\_CNTL\_REG}{0x0078}\label{regdesc:PMUHPSLEEPHPSYSCNTLREG}
\regfield{(reserved)}{2}{30}{00}%
\regfield{PMU\_HP\_SLEEP\_DIG\_CPU\_STALL}{1}{29}{{0}}%
\regfield{PMU\_HP\_SLEEP\_DIG\_PAUSE\_WDT}{1}{28}{{0}}%
\regfield{PMU\_HP\_SLEEP\_DIG\_PAD\_SLP\_SEL}{1}{27}{{0}}%
\regfield{PMU\_HP\_SLEEP\_HP\_PAD\_HOLD\_ALL}{1}{26}{{0}}%
\regfield{PMU\_HP\_SLEEP\_LP\_PAD\_HOLD\_ALL}{1}{25}{{0}}%
\regfield{PMU\_HP\_SLEEP\_UART\_WAKEUP\_EN}{1}{24}{{0}}%
\regfield{(reserved)}{24}{0}{000000000000000000000000}%
\reglabel{Reset}\regnewline%
\begin{regdesc}\begin{reglist}
\label{fielddesc:PMUHPSLEEPUARTWAKEUPEN}\item [PMU\_HP\_SLEEP\_UART\_WAKEUP\_EN] 配置是否在 HP\_SLEEP 状态下使能 UART 唤醒功能。 \\
0：关闭唤醒功能\\
1：使能唤醒功能\\
(R/W)
\label{fielddesc:PMUHPSLEEPLPPADHOLDALL}\item [PMU\_HP\_SLEEP\_LP\_PAD\_HOLD\_ALL] 配置是否在 HP\_SLEEP 状态下保持 LP GPIO 的配置。\\
0：不保持\\
1：保持\\
(R/W)
\label{fielddesc:PMUHPSLEEPHPPADHOLDALL}\item [PMU\_HP\_SLEEP\_HP\_PAD\_HOLD\_ALL] 配置是否在 HP\_SLEEP 状态下保持 GPIO 的配置。\\
0：不保持\\
1：保持\\
(R/W)
\label{fielddesc:PMUHPSLEEPDIGPADSLPSEL}\item [PMU\_HP\_SLEEP\_DIG\_PAD\_SLP\_SEL] 配置是否在 HP\_SLEEP 状态下使用 GPIO 的 Light-sleep 模式配置。\\
0：正常工作模式的配置\\
1：Light-sleep 模式配置，详见章节 \ref{mod:iomuxgpio} \textit{\nameref{mod:iomuxgpio}} > 小节 \ref{sec:iomuxgpio-pin-func-light-sleep} \textit{\nameref{sec:iomuxgpio-pin-func-light-sleep}}。\\
(R/W)
\label{fielddesc:PMUHPSLEEPDIGPAUSEWDT}\item [PMU\_HP\_SLEEP\_DIG\_PAUSE\_WDT] 配置是否在 HP\_SLEEP 状态下暂停看门狗。 \\
0：不暂停\\
1：暂停\\
(R/W)
\label{fielddesc:PMUHPSLEEPDIGCPUSTALL}\item [PMU\_HP\_SLEEP\_DIG\_CPU\_STALL] 配置是否在 HP\_SLEEP 状态下暂停 CPU。 \\
0：不暂停\\
1：暂停\\
(R/W)
\end{reglist}\end{regdesc}
\end{register}


\begin{register}{H}{PMU\_HP\_SLEEP\_HP\_CK\_POWER\_REG}{0x007C}\label{regdesc:PMUHPSLEEPHPCKPOWERREG}
\regfield{(reserved)}{1}{31}{0}%
\regfield{PMU\_HP\_SLEEP\_XPD\_PLL}{4}{27}{{0}}%
\regfield{PMU\_HP\_SLEEP\_XPD\_PLL\_I2C}{4}{23}{{0}}%
\regfield{PMU\_HP\_SLEEP\_I2C\_RETENTION}{1}{22}{{0}}%
\regfield{PMU\_HP\_SLEEP\_I2C\_ISO\_EN}{1}{21}{{0}}%
\regfield{(reserved)}{21}{0}{000000000000000000000}%
\reglabel{Reset}\regnewline%
\begin{regdesc}\begin{reglist}
\label{fielddesc:PMUHPSLEEPI2CISOEN}\item [PMU\_HP\_SLEEP\_I2C\_ISO\_EN]配置是否在 HP\_SLEEP 状态打开 ANALOG\_I2C。 \\
0：关闭\\
1：开启\\
(R/W)
\label{fielddesc:PMUHPSLEEPI2CRETENTION}\item [PMU\_HP\_SLEEP\_I2C\_RETENTION] 配置是否在 HP\_SLEEP 状态使能 ANALOG\_I2C 进入 retention 状态。 \\
0：不使能\\
1：使能\\
(R/W)
\label{fielddesc:PMUHPSLEEPXPDPLLI2C}\item [PMU\_HP\_SLEEP\_XPD\_PLL\_I2C] 配置是否在 HP\_SLEEP 状态开启 PLL I2C 组，该组寄存器为 bitmap，每个位控制一个 PLL I2C。 \\
0：关闭\\
1：开启\\
(R/W)
\label{fielddesc:PMUHPSLEEPXPDPLL}\item [PMU\_HP\_SLEEP\_XPD\_PLL] 配置是否在 HP\_SLEEP 状态下开启 HP 系统的模拟电源域，每个bit 控制一个电源域，bit 0 对应 CPLL\_CLK，bit 1 对应 SPLL\_CLK，bit 2 对应 Audio PLL，bit 3 对应 SDIO PLL。\\
0：关闭\\
1：开启\\
(R/W)
\end{reglist}\end{regdesc}
\end{register}


\begin{register}{H}{PMU\_HP\_SLEEP\_BIAS\_REG}{0x0080}\label{regdesc:PMUHPSLEEPBIASREG}
\regfield{(reserved)}{7}{25}{0000000}%
\regfield{PMU\_HP\_SLEEP\_DCM\_MODE}{2}{23}{{0}}%
\regfield{PMU\_HP\_SLEEP\_DCM\_VSET}{5}{18}{{20}}%
\regfield{(reserved)}{18}{0}{000000000000000000}%
\reglabel{Reset}\regnewline%
\begin{regdesc}\begin{reglist}
\label{fielddesc:PMUHPSLEEPDCMVSET}\item [PMU\_HP\_SLEEP\_DCM\_VSET]  调节 HP\_SLEEP 下 DCDC 的电压。(R/W) \\
\label{fielddesc:PMUHPSLEEPDCMMODE}\item [PMU\_HP\_SLEEP\_DCM\_MODE] 配置是否在 HP\_SLEEP 状态下使能 DCDC 的状态。\\
0：不使能 \\
1：使能 \\
2, 3：保留\\
(R/W)
\end{reglist}\end{regdesc}
\end{register}


\begin{register}{H}{PMU\_HP\_SLEEP\_BACKUP\_REG}{0x0084}\label{regdesc:PMUHPSLEEPBACKUPREG}
\regfield{PMU\_HP\_ACTIVE2SLEEP\_BACKUP\_EN}{1}{31}{{0}}%
\regfield{(reserved)}{2}{29}{00}%
\regfield{PMU\_HP\_ACTIVE2SLEEP\_BACKUP\_MODE}{3}{26}{{0}}%
\regfield{(reserved)}{6}{20}{000000}%
\regfield{PMU\_HP\_ACTIVE2SLEEP\_BACKUP\_CLK\_SEL}{2}{18}{{0}}%
\regfield{(reserved)}{18}{0}{00000000000000000}%
\reglabel{Reset}\regnewline%
\begin{regdesc}\begin{reglist}
\label{fielddesc:PMUHPACTIVE2SLEEPBACKUPCLKSEL}\item [PMU\_HP\_ACTIVE2SLEEP\_BACKUP\_CLK\_SEL] 配置备份模块从 HP\_ACTIVE 切换为 HP\_SLEEP 时的功能时钟源。\\
0：时钟源为 XTAL\\
1：时钟源为 PLL\_CLK\\
2：时钟源为 RC\_FAST\_CLK\\
3：无效配置\\
(R/W)
\label{fielddesc:PMUHPACTIVE2SLEEPBACKUPMODE}\item [PMU\_HP\_ACTIVE2SLEEP\_BACKUP\_MODE] 配置从 HP\_ACTIVE 切换为 HP\_SLEEP 时数据搬移的方向和链表指针。\\
最高位：\\
0：存储器到外设\\
1：外设到存储器\\
低两位：\\
0：PAU\_LINK\_ADDR\_0 \\
1：PAU\_LINK\_ADDR\_1 \\
2：PAU\_LINK\_ADDR\_2 \\
3：PAU\_LINK\_ADDR\_3 \\
(R/W)
\label{fielddesc:PMUHPACTIVE2SLEEPBACKUPEN}\item [PMU\_HP\_ACTIVE2SLEEP\_BACKUP\_EN] 配置是否在从 HP\_ACTIVE 切换为 HP\_SLEEP 时使能数据搬移。\\
0：禁能\\
1：使能\\
(R/W)
\end{reglist}\end{regdesc}
\end{register}



\begin{register}{H}{PMU\_HP\_SLEEP\_BACKUP\_CLK\_REG}{0x0088}\label{regdesc:PMUHPSLEEPBACKUPCLKREG}
\regfield{PMU\_HP\_SLEEP\_BACKUP\_ICG\_FUNC\_EN}{32}{0}{{0}}%
\reglabel{Reset}\regnewline%
\begin{regdesc}\begin{reglist}
\label{fielddesc:PMUHPSLEEPBACKUPICGFUNCEN}\item [PMU\_HP\_SLEEP\_BACKUP\_ICG\_FUNC\_EN] 配置目标状态为 HP\_SLEEP 时不同外设功能时钟的开关情况，不同位代表不同外设，详见表 \ref{tab:lowpowm-hp-sys-peri-clk}。\\
0：关闭\\
1：打开\\
(R/W)
\end{reglist}\end{regdesc}
\end{register}


\begin{register}{H}{PMU\_HP\_SLEEP\_SYSCLK\_REG}{0x008C}\label{regdesc:PMUHPSLEEPSYSCLKREG}
\regfield{PMU\_HP\_SLEEP\_DIG\_SYS\_CLK\_SEL}{2}{30}{{0}}%
\regfield{PMU\_HP\_SLEEP\_ICG\_SLP\_SEL}{1}{29}{{0}}%
\regfield{PMU\_HP\_SLEEP\_SYS\_CLK\_SLP\_SEL}{1}{28}{{0}}%
\regfield{PMU\_HP\_SLEEP\_ICG\_SYS\_CLOCK\_EN}{1}{27}{{0}}%
\regfield{(reserved)}{27}{0}{00000000000000000000000000}%
\reglabel{Reset}\regnewline%
\begin{regdesc}\begin{reglist}
\label{fielddesc:PMUHPSLEEPICGSYSCLOCKEN}\item [PMU\_HP\_SLEEP\_ICG\_SYS\_CLOCK\_EN] 配置是否在 HP\_SLEEP 状态下打开 ROOT\_CLK。\\
0：关闭\\
1：打开\\
(R/W)
\label{fielddesc:PMUHPSLEEPSYSCLKSLPSEL}\item [PMU\_HP\_SLEEP\_SYS\_CLK\_SLP\_SEL] 配置是否在 HP\_SLEEP 状态下由 PMU 控制时钟源。\\
0：由 PCR 寄存器控制\\
1：由 PMU 控制\\
(R/W)
\label{fielddesc:PMUHPSLEEPICGSLPSEL}\item [PMU\_HP\_SLEEP\_ICG\_SLP\_SEL] 配置是否在 HP\_SLEEP 状态下是否由 PMU 控制各时钟门控的开关。\\
0：由 PCR 寄存器控制\\
1：由 PMU 控制\\
(R/W)
\label{fielddesc:PMUHPSLEEPDIGSYSCLKSEL}\item [PMU\_HP\_SLEEP\_DIG\_SYS\_CLK\_SEL] 配置 HP\_SLEEP 状态下 ROOT\_CLK 的时钟源。\\
0：时钟源为 XTAL\\
1：时钟源为 PLL\_CLK\\
2：时钟源为 RC\_FAST\_CLK\\
3：无效配置\\
(R/W)
\end{reglist}\end{regdesc}
\end{register}


\begin{register}{H}{PMU\_HP\_SLEEP\_HP\_REGULATOR0\_REG}{0x0090}\label{regdesc:PMUHPSLEEPHPREGULATOR0REG} \regfield{PMU\_HP\_SLEEP\_HP\_REGULATOR\_DBIAS}{5}{27}{{24}}%
\regfield{(reserved)}{4}{23}{0000}
 \regfield{PMU\_HP\_SLEEP\_REGULATOR\_VO1\_XPD}{1}{22}{{1}}%
\regfield{(reserved)}{3}{19}{000}
\regfield{PMU\_HP\_SLEEP\_HP\_REGULATOR\_XPD}{1}{18}{{1}}%
\regfield{(reserved)}{18}{0}{000000000000000000}%
\reglabel{Reset}\regnewline%
\begin{regdesc}\begin{reglist}

\label{fielddesc:PMUHPSLEEPHPREGULATORXPD}\item [PMU\_HP\_SLEEP\_HP\_REGULATOR\_XPD] 配置是否在 HP\_SLEEP 状态下打开 HP 系统调压器。\\
0：关闭 HP 系统调压器 \\
1：打开 HP 系统调压器\\
(R/W)
\label{fielddesc:PMUHPSLEEPREGULATORVO1XPD}\item[PMU\_HP\_SLEEP\_REGULATOR\_VO1\_XPD] 配置 HP\_SLEEP 状态下 VO1 的开关。\\
0：关闭 \\
1：打开 \\
(R/W)
\label{fielddesc:PMUHPSLEEPHPREGULATORDBIAS}\item [PMU\_HP\_SLEEP\_HP\_REGULATOR\_DBIAS] 在 HP\_SLEEP 状态下调节 HP 系统调压器的电压，寄存器值越大，电压越高。(R/W)
\end{reglist}\end{regdesc}
\end{register}




\begin{register}{H}{PMU\_HP\_SLEEP\_XTAL\_REG}{0x0098}\label{regdesc:PMUHPSLEEPXTALREG}
\regfield{PMU\_HP\_SLEEP\_XPD\_XTAL}{1}{31}{{1}}%
\regfield{(reserved)}{31}{0}{0000000000000000000000000000000}%
\reglabel{Reset}\regnewline%
\begin{regdesc}\begin{reglist}
\label{fielddesc:PMUHPSLEEPXPDXTAL}\item [PMU\_HP\_SLEEP\_XPD\_XTAL] 配置是否在 HP\_SLEEP 状态下打开 XTAL\_CLK 的模拟源。\\
0：关闭\\
1：打开\\
(R/W)
\end{reglist}\end{regdesc}
\end{register}

\begin{register}{H}{PMU\_HP\_SLEEP\_LP\_REGULATOR0\_REG}{0x009C}\label{regdesc:PMUHPSLEEPLPREGULATOR0REG}
\regfield{PMU\_HP\_SLEEP\_LP\_REGULATOR\_DBIAS}{5}{27}{{24}}%
\regfield{(reserved)}{4}{23}{0000}%
\regfield{PMU\_HP\_SLEEP\_LP\_REGULATOR\_XPD}{1}{22}{{1}}%
\regfield{(reserved)}{22}{0}{0000000000000000000000}%
\reglabel{Reset}\regnewline%
\begin{regdesc}\begin{reglist}
\label{fielddesc:PMUHPSLEEPLPREGULATORXPD}\item [PMU\_HP\_SLEEP\_LP\_REGULATOR\_XPD] 配置是否在 HP\_ACTIVE/HP\_SLEEP 状态下打开 HP 系统调压器。\\
0：关闭 LP 系统调压器 \\
1：打开 LP 系统调压器\\
(R/W)
\label{fielddesc:PMUHPSLEEPLPREGULATORDBIAS}\item [PMU\_HP\_SLEEP\_LP\_REGULATOR\_DBIAS] 在 HP\_ACTIVE/HP\_SLEEP 状态下调节 LP 调压器的电压，寄存器值越大，电压越高。(R/W)
\end{reglist}\end{regdesc}
\end{register}


\begin{register}{H}{PMU\_HP\_SLEEP\_LP\_DIG\_POWER\_REG}{0x00A8}\label{regdesc:PMUHPSLEEPLPDIGPOWERREG}
\regfield{PMU\_HP\_SLEEP\_PD\_LP\_PERI\_PD\_EN}{1}{31}{{0}}%
\regfield{PMU\_HP\_SLEEP\_LP\_MEM\_DSLP}{1}{30}{{0}}%
\regfield{PMU\_HP\_SLEEP\_VDDBAT\_MODE}{2}{28}{{0}}%
\regfield{PMU\_HP\_SLEEP\_BOD\_SOURCE\_SEL}{1}{27}{{0}}%
\regfield{PMU\_HP\_SLEEP\_LP\_PAD\_SLP\_SEL}{1}{26}{{0}}%
\regfield{(reserved)}{26}{0}{00000000000000000000000000}%
\reglabel{Reset}\regnewline%
\begin{regdesc}\begin{reglist}
\label{fielddesc:PMUHPSLEEPLPPADSLPSEL}\item [PMU\_HP\_SLEEP\_LP\_PAD\_SLP\_SEL] 配置是否在 HP\_ACTIVE/HP\_SLEEP 状态下使能 LP GPIO 功能。 \\
0：不使能\\
1：使能 \\
(R/W)
\label{fielddesc:PMUHPSLEEPBODSOURCESEL}\item [PMU\_HP\_SLEEP\_BOD\_SOURCE\_SEL] 配置在 HP\_ACTIVE/HP\_SLEEP 状态下使能欠压检测器。\\ (R/W)
\label{fielddesc:PMUHPSLEEPVDDBATMODE}\item [PMU\_HP\_SLEEP\_VDDBAT\_MODE] 配置是否在 HP\_ACTIVE/HP\_SLEEP 状态下使能 VBAT 电源。 \\
0：不使能 \\
1：使能 \\
(R/W)   \label{fielddesc:PMUHPSLEEPLPMEMDSLP}\item [PMU\_HP\_SLEEP\_LP\_MEM\_DSLP] 配置是否在 HP\_ACTIVE/HP\_SLEEP 状态下使低功耗管理系统的 LP SRAM 进入 Memory Deep-sleep 模式。 \\
0：不使能\\
1：使能 \\
(R/W)
\label{fielddesc:PMUHPSLEEPPDLPPERIPDEN}\item [PMU\_HP\_SLEEP\_PD\_LP\_PERI\_PD\_EN] 配置是否在 HP\_ACTIVE/HP\_SLEEP 状态下关闭 LP PD Peripherals 电源域。 \\
0：打开\\
1：关闭 \\
(R/W)
\end{reglist}\end{regdesc}
\end{register}


\begin{register}{H}{PMU\_HP\_SLEEP\_LP\_CK\_POWER\_REG}{0x00AC}\label{regdesc:PMUHPSLEEPLPCKPOWERREG}
\regfield{PMU\_HP\_SLEEP\_PD\_OSC\_CLK}{1}{31}{{0}}%
\regfield{PMU\_HP\_SLEEP\_XPD\_FOSC\_CLK}{1}{30}{{1}}%
\regfield{PMU\_HP\_SLEEP\_XPD\_RC32K}{1}{29}{{0}}%
\regfield{PMU\_HP\_SLEEP\_XPD\_XTAL32K}{1}{28}{{0}}%
\regfield{PMU\_HP\_SLEEP\_XPD\_LPPLL}{1}{27}{{0}}%
\regfield{(reserved)}{27}{0}{000000000000000000000000000}%
\reglabel{Reset}\regnewline%
\begin{regdesc}\begin{reglist}
\label{fielddesc:PMUHPSLEEPXPDLPPLL}\item [PMU\_HP\_SLEEP\_XPD\_LPPLL] 配置是否在 HP\_ACTIVE/HP\_SLEEP 状态下打开 PLL\_LP\_CLK。\\
0：关闭\\
1：打开\\
(R/W)
\label{fielddesc:PMUHPSLEEPXPDXTAL32K}\item [PMU\_HP\_SLEEP\_XPD\_XTAL32K] 配置是否在 HP\_ACTIVE/HP\_SLEEP 状态下打开 XTAL32K\_CLK。\\
0：关闭\\
1：打开\\
(R/W)
\label{fielddesc:PMUHPSLEEPXPDRC32K}\item [PMU\_HP\_SLEEP\_XPD\_RC32K] 配置是否在 HP\_ACTIVE/HP\_SLEEP 状态下打开 RC32K\_CLK。\\
0：关闭\\
1：打开\\
(R/W)
\label{fielddesc:PMUHPSLEEPXPDFOSCCLK}\item [PMU\_HP\_SLEEP\_XPD\_FOSC\_CLK] 配置是否在 HP\_ACTIVE/HP\_SLEEP 状态下打开 RC\_FAST\_CLK。\\
0：关闭\\
1：打开\\
(R/W)
\label{fielddesc:PMUHPSLEEPPDOSCCLK}\item [PMU\_HP\_SLEEP\_PD\_OSC\_CLK] 配置是否在 HP\_ACTIVE/HP\_SLEEP 状态下打开 RC\_SLOW\_CLK。\\
0：关闭\\
1：打开\\
(R/W)
\end{reglist}\end{regdesc}
\end{register}


\begin{register}{H}{PMU\_LP\_SLEEP\_LP\_REGULATOR0\_REG}{0x00B4}\label{regdesc:PMULPSLEEPLPREGULATOR0REG}
\regfield{PMU\_LP\_SLEEP\_LP\_REGULATOR\_DBIAS}{5}{27}{{24}}%
\regfield{(reserved)}{4}{23}{0000}%
\regfield{PMU\_LP\_SLEEP\_LP\_REGULATOR\_XPD}{1}{22}{{1}}%
\regfield{(reserved)}{22}{0}{0000000000000000000000}%
\reglabel{Reset}\regnewline%
\begin{regdesc}\begin{reglist}
\label{fielddesc:PMULPSLEEPLPREGULATORXPD}\item [PMU\_LP\_SLEEP\_LP\_REGULATOR\_XPD] 配置是否在 LP\_SLEEP 状态下使能 LP 系统调压器。 \\
0：不使能\\
1：使能 \\
(R/W)
\label{fielddesc:PMULPSLEEPLPREGULATORDBIAS}\item [PMU\_LP\_SLEEP\_LP\_REGULATOR\_DBIAS] 在 LP\_SLEEP 状态下调节 LP 系统调压器的电压，寄存器值越大，电压越高。(R/W)
\end{reglist}\end{regdesc}
\end{register}




\begin{register}{H}{PMU\_LP\_SLEEP\_XTAL\_REG}{0x00BC}\label{regdesc:PMULPSLEEPXTALREG}
\regfield{PMU\_LP\_SLEEP\_XPD\_XTAL}{1}{31}{{1}}%
\regfield{(reserved)}{31}{0}{0000000000000000000000000000000}%
\reglabel{Reset}\regnewline%
\begin{regdesc}\begin{reglist}
\label{fielddesc:PMULPSLEEPXPDXTAL}\item [PMU\_LP\_SLEEP\_XPD\_XTAL] 配置是否在 LP\_SLEEP 状态下打开 XTAL\_CLK 的模拟源。\\
0：关闭\\
1：打开\\
(R/W)
\end{reglist}\end{regdesc}
\end{register}


\begin{register}{H}{PMU\_LP\_SLEEP\_LP\_DIG\_POWER\_REG}{0x00C0}\label{regdesc:PMULPSLEEPLPDIGPOWERREG}
\regfield{PMU\_LP\_SLEEP\_PD\_LP\_PERI\_PD\_EN}{1}{31}{{0}}%
\regfield{PMU\_LP\_SLEEP\_LP\_MEM\_DSLP}{1}{30}{{0}}%
\regfield{PMU\_LP\_SLEEP\_VDDBAT\_MODE}{2}{28}{{0}}%
\regfield{PMU\_LP\_SLEEP\_BOD\_SOURCE\_SEL}{1}{27}{{0}}%
\regfield{PMU\_LP\_SLEEP\_LP\_PAD\_SLP\_SEL}{1}{26}{{0}}%
\regfield{(reserved)}{26}{0}{00000000000000000000000000}%
\reglabel{Reset}\regnewline%
\begin{regdesc}\begin{reglist}
\label{fielddesc:PMULPSLEEPLPPADSLPSEL}\item [PMU\_LP\_SLEEP\_LP\_PAD\_SLP\_SEL] 配置是否在 LP\_SLEEP 状态下使能 LP GPIO 功能。 \\
0：不使能\\
1：使能 \\
(R/W)
\label{fielddesc:PMULPSLEEPBODSOURCESEL}\item [PMU\_LP\_SLEEP\_BOD\_SOURCE\_SEL] 配置在 LP\_SLEEP 状态下使能欠压检测器。\\ (R/W)
\label{fielddesc:PMULPSLEEPVDDBATMODE}\item [PMU\_LP\_SLEEP\_VDDBAT\_MODE] 配置是否在 LP\_SLEEP 状态下使能 VBAT 电源。 \\
0：不使能\\
1：使能\\
(R/W)
\label{fielddesc:PMULPSLEEPLPMEMDSLP}\item [PMU\_LP\_SLEEP\_LP\_MEM\_DSLP] 配置是否在 LP\_SLEEP 状态下使低功耗管理系统的 LP SRAM 进入 Memory Deep-sleep 模式。 \\
0：不使能\\
1：使能 \\
(R/W)
\label{fielddesc:PMULPSLEEPPDLPPERIPDEN}\item [PMU\_LP\_SLEEP\_PD\_LP\_PERI\_PD\_EN] 配置是否在 LP\_SLEEP 状态下关闭 LP PD Peripherals 电源域。 \\
0：打开\\
1：关闭 \\
(R/W)
\end{reglist}\end{regdesc}
\end{register}

\begin{register}{H}{PMU\_LP\_SLEEP\_LP\_CK\_POWER\_REG}{0x00C4}\label{regdesc:PMULPSLEEPLPCKPOWERREG}
\regfield{PMU\_LP\_SLEEP\_PD\_OSC\_CLK}{1}{31}{{0}}%
\regfield{PMU\_LP\_SLEEP\_XPD\_FOSC\_CLK}{1}{30}{{1}}%
\regfield{PMU\_LP\_SLEEP\_XPD\_RC32K}{1}{29}{{0}}%
\regfield{PMU\_LP\_SLEEP\_XPD\_XTAL32K}{1}{28}{{0}}%
\regfield{PMU\_LP\_SLEEP\_XPD\_LPPLL}{1}{27}{{0}}%
\regfield{(reserved)}{27}{0}{000000000000000000000000000}%
\reglabel{Reset}\regnewline%
\begin{regdesc}\begin{reglist}
\label{fielddesc:PMULPSLEEPXPDLPPLL}\item [PMU\_LP\_SLEEP\_XPD\_LPPLL] 配置是否在 LP\_SLEEP 状态下打开 PLL\_LP\_CLK。\\
0：关闭\\
1：打开\\
(R/W)
\label{fielddesc:PMULPSLEEPXPDXTAL32K}\item [PMU\_LP\_SLEEP\_XPD\_XTAL32K] 配置是否在 LP\_SLEEP 状态下打开 XTAL32K\_CLK。\\
0：关闭\\
1：打开\\
(R/W)
\label{fielddesc:PMULPSLEEPXPDRC32K}\item [PMU\_LP\_SLEEP\_XPD\_RC32K] 配置是否在 LP\_SLEEP 状态下打开 RC32K\_CLK。\\
0：关闭\\
1：打开\\
(R/W)
\label{fielddesc:PMULPSLEEPXPDFOSCCLK}\item [PMU\_LP\_SLEEP\_XPD\_FOSC\_CLK] 配置是否在 LP\_SLEEP 状态下打开 RC\_FAST\_CLK。\\
0：关闭\\
1：打开\\
(R/W)
\label{fielddesc:PMULPSLEEPPDOSCCLK}\item [PMU\_LP\_SLEEP\_PD\_OSC\_CLK] 配置是否在 LP\_SLEEP 状态下关闭 RC\_SLOW\_CLK。\\
0：打开\\
1：关闭\\
(R/W)
\end{reglist}\end{regdesc}
\end{register}


\begin{register}{H}{PMU\_IMM\_HP\_CK\_POWER\_REG}{0x00CC}\label{regdesc:PMUIMMHPCKPOWERREG}
\regfield{PMU\_TIE\_HIGH\_XPD\_XTAL}{1}{31}{{0}}%
\regfield{PMU\_TIE\_HIGH\_XPD\_PLL}{4}{27}{{0}}%
\regfield{PMU\_TIE\_HIGH\_XPD\_PLL\_I2C}{4}{23}{{0}}%
\regfield{PMU\_TIE\_HIGH\_I2C\_RETENTION}{1}{22}{{0}}%
\regfield{PMU\_TIE\_HIGH\_GLOBAL\_XTAL\_ICG}{1}{21}{{0}}%
\regfield{PMU\_TIE\_HIGH\_GLOBAL\_PLL\_ICG}{4}{17}{{0}}%
\regfield{reserved}{1}{16}{{0}}%
\regfield{PMU\_TIE\_LOW\_XPD\_XTAL}{1}{15}{{0}}%
\regfield{PMU\_TIE\_LOW\_XPD\_PLL}{4}{11}{{0}}%
\regfield{PMU\_TIE\_LOW\_XPD\_PLL\_I2C}{4}{7}{{0}}%
\regfield{PMU\_TIE\_LOW\_I2C\_RETENTION}{1}{6}{{0}}%
\regfield{PMU\_TIE\_LOW\_GLOBAL\_XTAL\_ICG}{1}{5}{{0}}%
\regfield{PMU\_TIE\_LOW\_GLOBAL\_PLL\_ICG}{4}{1}{{0}}%
\regfield{reserved}{1}{0}{{0}}%
\reglabel{Reset}\regnewline%
\begin{regdesc}\begin{reglist}
\label{fielddesc:PMUTIELOWGLOBALPLLICG}\item [PMU\_TIE\_LOW\_GLOBAL\_PLL\_ICG] 配置关闭全局 PLL 组的时钟门控位图。\\
0：无操作 \\
1：关闭  \\
(WT)
\label{fielddesc:PMUTIELOWGLOBALXTALICG}\item [PMU\_TIE\_LOW\_GLOBAL\_XTAL\_ICG] 配置关闭全局 XTAL 的时钟门控。\\
0：无操作 \\
1：关闭  \\
(WT)
\label{fielddesc:PMUTIELOWI2CRETENTION}\item [PMU\_TIE\_LOW\_I2C\_RETENTION] 配置关闭 ANALOG I2C retention。\\
0：无操作 \\
1：关闭  \\
(WT)
\label{fielddesc:PMUTIELOWXPDPLLI2C}\item [PMU\_TIE\_LOW\_XPD\_PLL\_I2C] 配置关闭全局 PLL I2C 组的时钟电源位图。\\
0：无操作 \\
1：关闭  \\
(WT)
\label{fielddesc:PMUTIELOWXPDPLL}\item [PMU\_TIE\_LOW\_XPD\_PLL] 配置关闭全局 PLL  组的时钟电源位图。\\
0：无操作 \\
1：关闭  \\
(WT)
\label{fielddesc:PMUTIELOWXPDXTAL}\item [PMU\_TIE\_LOW\_XPD\_XTAL]  配置关闭全局 XTAL 的时钟电源。\\
0：无操作 \\
1：关闭  \\
(WT)
\label{fielddesc:PMUTIEHIGHGLOBALPLLICG}\item [PMU\_TIE\_HIGH\_GLOBAL\_PLL\_ICG] 配置开启全局 PLL 组的时钟门控位图。\\
0：无操作 \\
1：开启  \\
(WT)

\item [见下页...]
\end{reglist}\end{regdesc}
\end{register}

\addtocounter{Regfloat}{-1}
\begin{register}{H}{PMU\_IMM\_HP\_CK\_POWER\_REG}{0x00CC}
\begin{regdesc}\begin{reglist}
\item [接上页...]

\label{fielddesc:PMUTIEHIGHGLOBALXTALICG}\item [PMU\_TIE\_HIGH\_GLOBAL\_XTAL\_ICG] 配置开启全局 XTAL 的时钟门控。\\
0：无操作 \\
1：开启  \\
(WT)
\label{fielddesc:PMUTIEHIGHI2CRETENTION}\item [PMU\_TIE\_HIGH\_I2C\_RETENTION] 配置开启 ANALOG I2C retention。\\
0：无操作 \\
1：开启  \\
(WT)
\label{fielddesc:PMUTIEHIGHXPDPLLI2C}\item [PMU\_TIE\_HIGH\_XPD\_PLL\_I2C]配置开启全局 PLL I2C 组的时钟电源位图。\\
0：无操作 \\
1：开启  \\
(WT)
\label{fielddesc:PMUTIEHIGHXPDPLL}\item [PMU\_TIE\_HIGH\_XPD\_PLL] 配置开启全局 PLL 组的时钟电源位图。\\
0：无操作 \\
1：开启  \\
(WT)
\label{fielddesc:PMUTIEHIGHXPDXTAL}\item [PMU\_TIE\_HIGH\_XPD\_XTAL]  配置开启全局 XTAL 的时钟电源。\\
0：无操作 \\
1：开启  \\
(WT)
\end{reglist}\end{regdesc}
\end{register}


\begin{register}{H}{PMU\_IMM\_SLEEP\_SYSCLK\_REG}{0x00D0}\label{regdesc:PMUIMMSLEEPSYSCLKREG}
\regfield{PMU\_UPDATE\_DIG\_SYS\_CLK\_SEL}{1}{31}{{0}}%
\regfield{PMU\_TIE\_HIGH\_ICG\_SLP\_SEL}{1}{30}{{0}}%
\regfield{PMU\_TIE\_LOW\_ICG\_SLP\_SEL}{1}{29}{{0}}%
\regfield{PMU\_UPDATE\_DIG\_ICG\_SWITCH}{1}{28}{{0}}%
\regfield{(reserved)}{28}{0}{0000000000000000000000000000}%
\reglabel{Reset}\regnewline%
\begin{regdesc}\begin{reglist}
\label{fielddesc:PMUUPDATEDIGICGSWITCH}\item [PMU\_UPDATE\_DIG\_ICG\_SWITCH] 立即更新功能时钟门控状态切换。(WT)
\label{fielddesc:PMUTIELOWICGSLPSEL}\item [PMU\_TIE\_LOW\_ICG\_SLP\_SEL] 配置关闭功能模块时钟门控的 PMU 控制。\\
0：无操作 \\
1：关闭  \\
(WT)
\label{fielddesc:PMUTIEHIGHICGSLPSEL}\item [PMU\_TIE\_HIGH\_ICG\_SLP\_SEL] 配置开启功能模块时钟门控的 PMU 控制。\\
0：无操作 \\
1：开启  \\
(WT)
\label{fielddesc:PMUUPDATEDIGSYSCLKSEL}\item [PMU\_UPDATE\_DIG\_SYS\_CLK\_SEL] 立即更新系统时钟选择。\\
0：无操作\\
1：更新当前 PMU 状态的 SYS\_CLK\_SEL 配置到使用端\\
(WT)
\end{reglist}\end{regdesc}
\end{register}


\begin{register}{H}{PMU\_IMM\_HP\_FUNC\_ICG\_REG}{0x00D4}\label{regdesc:PMUIMMHPFUNCICGREG}
\regfield{PMU\_UPDATE\_DIG\_ICG\_FUNC\_EN}{1}{31}{{0}}%
\regfield{(reserved)}{31}{0}{0000000000000000000000000000000}%
\reglabel{Reset}\regnewline%
\begin{regdesc}\begin{reglist}
\label{fielddesc:PMUUPDATEDIGICGFUNCEN}\item [PMU\_UPDATE\_DIG\_ICG\_FUNC\_EN] 立即更新功能模块时钟门控的位图。(WT)
\end{reglist}\end{regdesc}
\end{register}


\begin{register}{H}{PMU\_IMM\_HP\_APB\_ICG\_REG}{0x00D8}\label{regdesc:PMUIMMHPAPBICGREG}
\regfield{PMU\_UPDATE\_DIG\_ICG\_APB\_EN}{1}{31}{{0}}%
\regfield{(reserved)}{31}{0}{0000000000000000000000000000000}%
\reglabel{Reset}\regnewline%
\begin{regdesc}\begin{reglist}
\label{fielddesc:PMUUPDATEDIGICGAPBEN}\item [PMU\_UPDATE\_DIG\_ICG\_APB\_EN]  立即更新 APB 时钟门控的位图。(WT)
\end{reglist}\end{regdesc}
\end{register}






\begin{register}{H}{PMU\_IMM\_PAD\_HOLD\_ALL\_REG}{0x00E4}\label{regdesc:PMUIMMPADHOLDALLREG}
\regfield{PMU\_TIE\_LOW\_HP\_PAD\_HOLD\_ALL}{1}{31}{{0}}%
\regfield{PMU\_TIE\_HIGH\_HP\_PAD\_HOLD\_ALL}{1}{30}{{0}}%
\regfield{PMU\_TIE\_LOW\_LP\_PAD\_HOLD\_ALL}{1}{29}{{0}}%
\regfield{PMU\_TIE\_HIGH\_LP\_PAD\_HOLD\_ALL}{1}{28}{{0}}%
\regfield{PMU\_TIE\_LOW\_PAD\_SLP\_SEL}{1}{27}{{0}}%
\regfield{PMU\_TIE\_HIGH\_PAD\_SLP\_SEL}{1}{26}{{0}}%
\regfield{(reserved)}{23}{3}{00000000000000000000000}%
\regfield{PMU\_HP\_PAD\_HOLD\_ALL}{1}{2}{{0}}%
\regfield{PMU\_LP\_PAD\_HOLD\_ALL}{1}{1}{{0}}%
\regfield{PMU\_PAD\_SLP\_SEL}{1}{0}{{0}}%
\reglabel{Reset}\regnewline%
\begin{regdesc}\begin{reglist}
\label{fielddesc:PMUPADSLPSEL}\item [PMU\_PAD\_SLP\_SEL] 表示当前 GPIO 的功能是否使用睡眠的配置。\\
0：不是 \\
1：是  \\
(RO)
\label{fielddesc:PMULPPADHOLDALL}\item [PMU\_LP\_PAD\_HOLD\_ALL] 表示当前 LP PAD 是否处于 HOLD 状态。\\
0：不是 \\
1：是  \\
(RO)
\label{fielddesc:PMUHPPADHOLDALL}\item [PMU\_HP\_PAD\_HOLD\_ALL] 表示当前 HP PAD 是否处于 HOLD 状态。\\
0：不是 \\
1：是  \\
(RO)
\label{fielddesc:PMUTIEHIGHPADSLPSEL}\item [PMU\_TIE\_HIGH\_PAD\_SLP\_SEL] 立即配置 GPIO 功能使用睡眠配置。 (WT)
\label{fielddesc:PMUTIELOWPADSLPSEL}\item [PMU\_TIE\_LOW\_PAD\_SLP\_SEL] 立即解除 GPIO 功能使用睡眠配置。 (WT)
\label{fielddesc:PMUTIEHIGHLPPADHOLDALL}\item [PMU\_TIE\_HIGH\_LP\_PAD\_HOLD\_ALL] 立即配置 LP PAD 进入 HOLD 状态。(WT)
\label{fielddesc:PMUTIELOWLPPADHOLDALL}\item [PMU\_TIE\_LOW\_LP\_PAD\_HOLD\_ALL] 立即解除 LP PAD 进入 HOLD 状态。(WT)
\label{fielddesc:PMUTIEHIGHHPPADHOLDALL}\item [PMU\_TIE\_HIGH\_HP\_PAD\_HOLD\_ALL] 立即配置 HP PAD 进入 HOLD 状态。(WT)
\label{fielddesc:PMUTIELOWHPPADHOLDALL}\item [PMU\_TIE\_LOW\_HP\_PAD\_HOLD\_ALL] 立即解除 HP PAD 进入 HOLD 状态。 (WT)
\end{reglist}\end{regdesc}
\end{register}




























\begin{register}{H}{PMU\_POWER\_CK\_WAIT\_CNTL\_REG}{0x011C}\label{regdesc:PMUPOWERCKWAITCNTLREG}
\regfield{PMU\_WAIT\_PLL\_STABLE}{16}{16}{{256}}%
\regfield{PMU\_WAIT\_XTL\_STABLE}{16}{0}{{256}}%
\reglabel{Reset}\regnewline%
\begin{regdesc}\begin{reglist}
\label{fielddesc:PMUWAITXTLSTABLE}\item [PMU\_WAIT\_XTL\_STABLE] 配置开启 XTAL 电源后，开启 XTAL 全局时钟门控的等待时间。单位是 LP\_DYN\_FAST\_CLK。(R/W)
\label{fielddesc:PMUWAITPLLSTABLE}\item [PMU\_WAIT\_PLL\_STABLE] 配置开启 PLL 电源后，开启 PLL 全局时钟门控的等待时间。单位是 LP\_DYN\_FAST\_CLK。(R/W)
\end{reglist}\end{regdesc}
\end{register}


\begin{register}{H}{PMU\_SLP\_WAKEUP\_CNTL0\_REG}{0x0120}\label{regdesc:PMUSLPWAKEUPCNTL0REG}
\regfield{PMU\_SLEEP\_REQ}{1}{31}{{0}}%
\regfield{(reserved)}{31}{0}{0000000000000000000000000000000}%
\reglabel{Reset}\regnewline%
\begin{regdesc}\begin{reglist}
\label{fielddesc:PMUSLEEPREQ}\item [PMU\_SLEEP\_REQ] 配置是否将芯片的 PMU 状态切换至 HP\_SLEEP 或 LP\_SLEEP。\\
0：不切换\\
1：切换，目标状态取决于 LP CPU 的工作状态。\\
(WT)
\end{reglist}\end{regdesc}
\end{register}


\begin{register}{H}{PMU\_SLP\_WAKEUP\_CNTL1\_REG}{0x0124}\label{regdesc:PMUSLPWAKEUPCNTL1REG}
\regfield{PMU\_SLP\_REJECT\_EN}{1}{31}{{0}}%
\regfield{PMU\_SLEEP\_REJECT\_ENA}{31}{0}{{0}}%
\reglabel{Reset}\regnewline%
\begin{regdesc}\begin{reglist}
\label{fielddesc:PMUSLEEPREJECTENA}\item [PMU\_SLEEP\_REJECT\_ENA] 配置拒绝睡眠信号源。配置值与拒绝睡眠信号源的对应关系见表 \ref{tab:lowpowm-wakeup-source}。(R/W)
\label{fielddesc:PMUSLPREJECTEN}\item [PMU\_SLP\_REJECT\_EN] 配置是否使能拒绝睡眠功能。\\
0：禁能\\
1：使能\\
(R/W)
\end{reglist}\end{regdesc}
\end{register}


\begin{register}{H}{PMU\_SLP\_WAKEUP\_CNTL2\_REG}{0x0128}\label{regdesc:PMUSLPWAKEUPCNTL2REG}
\regfield{(reserved)}{1}{31}{0}%
\regfield{PMU\_WAKEUP\_ENA}{31}{0}{{0}}%
\reglabel{Reset}\regnewline%
\begin{regdesc}\begin{reglist}
\label{fielddesc:PMUWAKEUPENA}\item [PMU\_WAKEUP\_ENA] 配置唤醒源。配置值与唤醒源的对应关系见表 \ref{tab:lowpowm-wakeup-source}。\\ (R/W)
\end{reglist}\end{regdesc}
\end{register}


\begin{register}{H}{PMU\_SLP\_WAKEUP\_CNTL3\_REG}{0x012C}\label{regdesc:PMUSLPWAKEUPCNTL3REG}
\regfield{(reserved)}{14}{18}{00000000000000}%
\regfield{PMU\_SLEEP\_PRT\_SEL}{2}{16}{{0}}%
\regfield{PMU\_HP\_MIN\_SLP\_VAL}{8}{8}{{0}}%
\regfield{PMU\_LP\_MIN\_SLP\_VAL}{8}{0}{{0}}%
\reglabel{Reset}\regnewline%
\begin{regdesc}\begin{reglist}
\label{fielddesc:PMULPMINSLPVAL}\item [PMU\_LP\_MIN\_SLP\_VAL] 配置进入 LP\_SLEEP 的最小睡眠时间。单位是 LP\_DYN\_SLOW\_CLK。(R/W)
\label{fielddesc:PMUHPMINSLPVAL}\item [PMU\_HP\_MIN\_SLP\_VAL] 配置进入 HP\_SLEEP 的最小睡眠时间。单位是 LP\_DYN\_SLOW\_CLK。(R/W)
\label{fielddesc:PMUSLEEPPRTSEL}\item [PMU\_SLEEP\_PRT\_SEL] 配置最小睡眠时间模式。\\
0, 1, 3：保留\\
2：同时保护 LP\_SLEEP 和 HP\_SLEEP \\
(R/W)
\end{reglist}\end{regdesc}
\end{register}


\begin{register}{H}{PMU\_SLP\_WAKEUP\_CNTL4\_REG}{0x0130}\label{regdesc:PMUSLPWAKEUPCNTL4REG}
\regfield{PMU\_SLP\_REJECT\_CAUSE\_CLR}{1}{31}{{0}}%
\regfield{(reserved)}{31}{0}{0000000000000000000000000000000}%
\reglabel{Reset}\regnewline%
\begin{regdesc}\begin{reglist}
\label{fielddesc:PMUSLPREJECTCAUSECLR}\item [PMU\_SLP\_REJECT\_CAUSE\_CLR] 写 1 清除 \hyperref[fielddesc:PMUREJECTCAUSE]{PMU\_REJECT\_CAUSE}。(WT)
\end{reglist}\end{regdesc}
\end{register}










\begin{register}{H}{PMU\_SLP\_WAKEUP\_STATUS0\_REG}{0x0144}\label{regdesc:PMUSLPWAKEUPSTATUS0REG}
\regfield{(reserved)}{1}{31}{0}%
\regfield{PMU\_WAKEUP\_CAUSE}{31}{0}{{0}}%
\reglabel{Reset}\regnewline%
\begin{regdesc}\begin{reglist}
\label{fielddesc:PMUWAKEUPCAUSE}\item [PMU\_WAKEUP\_CAUSE] 指示唤醒源。配置值与唤醒源的对应关系请见表 \ref{tab:lowpowm-wakeup-source}。(RO)
\end{reglist}\end{regdesc}
\end{register}


\begin{register}{H}{PMU\_SLP\_WAKEUP\_STATUS1\_REG}{0x0148}\label{regdesc:PMUSLPWAKEUPSTATUS1REG}
\regfield{(reserved)}{1}{31}{0}%
\regfield{PMU\_REJECT\_CAUSE}{31}{0}{{0}}%
\reglabel{Reset}\regnewline%
\begin{regdesc}\begin{reglist}
\label{fielddesc:PMUREJECTCAUSE}\item [PMU\_REJECT\_CAUSE] 指示拒绝唤醒的事件源。配置值与唤醒源的对应关系请见表 \ref{tab:lowpowm-wakeup-source}。(RO)
\end{reglist}\end{regdesc}
\end{register}


\begin{register}{H}{PMU\_HP\_CK\_CNTL\_REG}{0x0154}\label{regdesc:PMUHPCKCNTLREG}
\regfield{(reserved)}{16}{16}{0000000000000000}%
\regfield{PMU\_SWITCH\_ICG\_CNTL\_WAIT}{8}{8}{{10}}%
\regfield{PMU\_MODIFY\_ICG\_CNTL\_WAIT}{8}{0}{{10}}%
\reglabel{Reset}\regnewline%
\begin{regdesc}\begin{reglist}
\label{fielddesc:PMUMODIFYICGCNTLWAIT}\item [PMU\_MODIFY\_ICG\_CNTL\_WAIT] 配置修改功能模块时钟门控的等待时间。单位是 LP\_DYN\_FAST\_CLK。(R/W)
\label{fielddesc:PMUSWITCHICGCNTLWAIT}\item [PMU\_SWITCH\_ICG\_CNTL\_WAIT]配置切换功能模块时钟门控的等待时间。单位是 LP\_DYN\_FAST\_CLK。(R/W)
\end{reglist}\end{regdesc}
\end{register}




\begin{register}{H}{PMU\_RF\_PWC\_REG}{0x015C}\label{regdesc:PMURFPWCREG}
\regfield{(reserved)}{4}{28}{0000}%
\regfield{PMU\_XPD\_PERIF\_I2C}{1}{27}{{1}}%
\regfield{PMU\_PERIF\_I2C\_RSTB}{1}{26}{{0}}%
\regfield{PMU\_SDIO\_PLL\_XPD}{1}{25}{{0}}%
\regfield{PMU\_MSPI\_PHY\_XPD}{1}{24}{{0}}%
\regfield{(reserved)}{24}{0}{000000000000000000000000}%
\reglabel{Reset}\regnewline%
\begin{regdesc}\begin{reglist}
\label{fielddesc:PMUMSPIPHYXPD}\item [PMU\_MSPI\_PHY\_XPD] 配置是否开启 MSPI PHY。 \\
0：关闭 \\
1：开启 \\
(R/W)
\label{fielddesc:PMUSDIOPLLXPD}\item [PMU\_SDIO\_PLL\_XPD] 配置是否开启 SDIO PLL。 \\
0：关闭 \\
1：开启 \\
(R/W)
\label{fielddesc:PMUPERIFI2CRSTB}\item [PMU\_PERIF\_I2C\_RSTB]  配置是否使能 PERI I2C 的复位。 \\
0：复位 \\
1：释放 \\
(R/W)
\label{fielddesc:PMUXPDPERIFI2C}\item [PMU\_XPD\_PERIF\_I2C] 配置是否开启 PERI I2C 电源。 \\
0：关闭 \\
1：开启 \\
(R/W)

\end{reglist}\end{regdesc}
\end{register}


\begin{register}{H}{PMU\_INT\_RAW\_REG}{0x0164}\label{regdesc:PMUINTRAWREG}
\regfield{PMU\_SOC\_WAKEUP\_INT\_RAW}{1}{31}{{0}}%
\regfield{PMU\_SOC\_SLEEP\_REJECT\_INT\_RAW}{1}{30}{{0}}%
\regfield{PMU\_SW\_INT\_RAW}{1}{29}{{0}}%
\regfield{PMU\_SDIO\_IDLE\_INT\_RAW}{1}{28}{{0}}%
\regfield{PMU\_LP\_CPU\_EXC\_INT\_RAW}{1}{27}{{0}}%
\regfield{(reserved)}{5}{22}{000000}%
\regfield{PMU\_0P2A\_CNT\_TARGET1\_REACH\_1\_HP\_INT\_RAW}{1}{21}{{0}}%
\regfield{PMU\_0P2A\_CNT\_TARGET0\_REACH\_1\_HP\_INT\_RAW}{1}{20}{{0}}%
\regfield{PMU\_0P2A\_CNT\_TARGET1\_REACH\_0\_HP\_INT\_RAW}{1}{19}{{0}}%
\regfield{PMU\_0P2A\_CNT\_TARGET0\_REACH\_0\_HP\_INT\_RAW}{1}{18}{{0}}%
\regfield{PMU\_0P1A\_CNT\_TARGET1\_REACH\_1\_HP\_INT\_RAW}{1}{17}{{0}}%
\regfield{PMU\_0P1A\_CNT\_TARGET0\_REACH\_1\_HP\_INT\_RAW}{1}{16}{{0}}%
\regfield{PMU\_0P1A\_CNT\_TARGET1\_REACH\_0\_HP\_INT\_RAW}{1}{15}{{0}}%
\regfield{PMU\_0P1A\_CNT\_TARGET0\_REACH\_0\_HP\_INT\_RAW}{1}{14}{{0}}%
\regfield{(reserved)}{14}{0}{00000000000000}%
\reglabel{Reset}\regnewline%
\begin{regdesc}\begin{reglist}
\label{fielddesc:PMU0P1ACNTTARGET0REACH0HPINTRAW}\item [PMU\_0P1A\_CNT\_TARGET0\_REACH\_0\_HP\_INT\_RAW] \hyperref[int:PMU0P1ACNTTARGET0REACH0INT]{PMU\_0P1A\_CNT\_TARGET0\_REACH\_0\_INT} 的原始中断状态。(R/WTC/SS)

\label{fielddesc:PMU0P1ACNTTARGET1REACH0HPINTRAW}\item [PMU\_0P1A\_CNT\_TARGET1\_REACH\_0\_HP\_INT\_RAW] \hyperref[int:PMU0P1ACNTTARGET1REACH0INT]{PMU\_0P1A\_CNT\_TARGET1\_REACH\_0\_INT} 的原始中断状态。(R/WTC/SS)
\label{fielddesc:PMU0P1ACNTTARGET0REACH1HPINTRAW}\item [PMU\_0P1A\_CNT\_TARGET0\_REACH\_1\_HP\_INT\_RAW] \hyperref[int:PMU0P1ACNTTARGET0REACH1INT]{PMU\_0P1A\_CNT\_TARGET0\_REACH\_1\_INT} 的原始中断状态。 (R/WTC/SS)
\label{fielddesc:PMU0P1ACNTTARGET1REACH1HPINTRAW}\item [PMU\_0P1A\_CNT\_TARGET1\_REACH\_1\_HP\_INT\_RAW] \hyperref[int:PMU0P1ACNTTARGET1REACH1INT]{PMU\_0P1A\_CNT\_TARGET1\_REACH\_1\_INT} 的原始中断状态。(R/WTC/SS)
\label{fielddesc:PMU0P2ACNTTARGET0REACH0HPINTRAW}\item [PMU\_0P2A\_CNT\_TARGET0\_REACH\_0\_HP\_INT\_RAW] \hyperref[int:PMU0P2ACNTTARGET0REACH0INT]{PMU\_0P2A\_CNT\_TARGET0\_REACH\_0\_INT} 的原始中断状态。(R/WTC/SS)
\label{fielddesc:PMU0P2ACNTTARGET1REACH0HPINTRAW}\item [PMU\_0P2A\_CNT\_TARGET1\_REACH\_0\_HP\_INT\_RAW] \hyperref[int:PMU0P2ACNTTARGET1REACH0INT]{PMU\_0P2A\_CNT\_TARGET1\_REACH\_0\_INT} 的原始中断状态。(R/WTC/SS)
\label{fielddesc:PMU0P2ACNTTARGET0REACH1HPINTRAW}\item [PMU\_0P2A\_CNT\_TARGET0\_REACH\_1\_HP\_INT\_RAW] \hyperref[int:PMU0P2ACNTTARGET0REACH1INT]{PMU\_0P2A\_CNT\_TARGET0\_REACH\_1\_INT} 的原始中断状态。(R/WTC/SS)
\label{fielddesc:PMU0P2ACNTTARGET1REACH1HPINTRAW}\item [PMU\_0P2A\_CNT\_TARGET1\_REACH\_1\_HP\_INT\_RAW] \hyperref[int:PMU0P2ACNTTARGET1REACH1INT]{PMU\_0P2A\_CNT\_TARGET1\_REACH\_1\_INT} 的原始中断状态。(R/WTC/SS)
\label{fielddesc:PMULPCPUEXCINTRAW}\item [PMU\_LP\_CPU\_EXC\_INT\_RAW] \hyperref[int:PMULPCPUEXCINT]{PMU\_LP\_CPU\_EXC\_INT} 的原始中断状态。(R/WTC/SS)
\label{fielddesc:PMUSDIOIDLEINTRAW}\item [PMU\_SDIO\_IDLE\_INT\_RAW] \hyperref[int:PMUSDIOIDLEINT]{PMU\_SDIO\_IDLE\_INT} 的原始中断状态。(R/WTC/SS)
\label{fielddesc:PMUSWINTRAW}\item [PMU\_SW\_INT\_RAW] \hyperref[int:PMUSWINT]{PMU\_SW\_INT} 的原始中断状态。 (R/WTC/SS)
\label{fielddesc:PMUSOCSLEEPREJECTINTRAW}\item [PMU\_SOC\_SLEEP\_REJECT\_INT\_RAW] \hyperref[int:PMUSOCSLEEPREJECTINT]{PMU\_SOC\_SLEEP\_REJECT\_INT} 的原始中断状态。(R/WTC/SS)
\label{fielddesc:PMUSOCWAKEUPINTRAW}\item [PMU\_SOC\_WAKEUP\_INT\_RAW] \hyperref[int:PMUSOCWAKEUPINT]{PMU\_SOC\_WAKEUP\_INT} 的原始中断状态。(R/WTC/SS)
\end{reglist}\end{regdesc}
\end{register}


\begin{register}{H}{PMU\_HP\_INT\_ST\_REG}{0x0168}\label{regdesc:PMUHPINTSTREG}
\regfield{PMU\_SOC\_WAKEUP\_INT\_ST}{1}{31}{{0}}%
\regfield{PMU\_SOC\_SLEEP\_REJECT\_INT\_ST}{1}{30}{{0}}%
\regfield{PMU\_SW\_INT\_ST}{1}{29}{{0}}%
\regfield{PMU\_SDIO\_IDLE\_INT\_ST}{1}{28}{{0}}%
\regfield{PMU\_LP\_CPU\_EXC\_INT\_ST}{1}{27}{{0}}%
\regfield{(reserved)}{5}{22}{000000}%
\regfield{PMU\_0P2A\_CNT\_TARGET1\_REACH\_1\_HP\_INT\_ST}{1}{21}{{0}}%
\regfield{PMU\_0P2A\_CNT\_TARGET0\_REACH\_1\_HP\_INT\_ST}{1}{20}{{0}}%
\regfield{PMU\_0P2A\_CNT\_TARGET1\_REACH\_0\_HP\_INT\_ST}{1}{19}{{0}}%
\regfield{PMU\_0P2A\_CNT\_TARGET0\_REACH\_0\_HP\_INT\_ST}{1}{18}{{0}}%
\regfield{PMU\_0P1A\_CNT\_TARGET1\_REACH\_1\_HP\_INT\_ST}{1}{17}{{0}}%
\regfield{PMU\_0P1A\_CNT\_TARGET0\_REACH\_1\_HP\_INT\_ST}{1}{16}{{0}}%
\regfield{PMU\_0P1A\_CNT\_TARGET1\_REACH\_0\_HP\_INT\_ST}{1}{15}{{0}}%
\regfield{PMU\_0P1A\_CNT\_TARGET0\_REACH\_0\_HP\_INT\_ST}{1}{14}{{0}}%
\regfield{(reserved)}{14}{0}{00000000000000}%
\reglabel{Reset}\regnewline%
\begin{regdesc}\begin{reglist}
\label{fielddesc:PMU0P1ACNTTARGET0REACH0HPINTST}\item [PMU\_0P1A\_CNT\_TARGET0\_REACH\_0\_HP\_INT\_ST] \hyperref[int:PMU0P1ACNTTARGET0REACH0INT]{PMU\_0P1A\_CNT\_TARGET0\_REACH\_0\_INT} 的屏蔽中断状态。(RO)
\label{fielddesc:PMU0P1ACNTTARGET1REACH0HPINTST}\item [PMU\_0P1A\_CNT\_TARGET1\_REACH\_0\_HP\_INT\_ST] \hyperref[int:PMU0P1ACNTTARGET1REACH0INT]{PMU\_0P1A\_CNT\_TARGET1\_REACH\_0\_INT} 的屏蔽中断状态。(RO)
\label{fielddesc:PMU0P1ACNTTARGET0REACH1HPINTST}\item [PMU\_0P1A\_CNT\_TARGET0\_REACH\_1\_HP\_INT\_ST] \hyperref[int:PMU0P1ACNTTARGET0REACH1INT]{PMU\_0P1A\_CNT\_TARGET0\_REACH\_1\_INT} 的屏蔽中断状态。(RO)
\label{fielddesc:PMU0P1ACNTTARGET1REACH1HPINTST}\item [PMU\_0P1A\_CNT\_TARGET1\_REACH\_1\_HP\_INT\_ST] \hyperref[int:PMU0P1ACNTTARGET1REACH1INT]{PMU\_0P1A\_CNT\_TARGET1\_REACH\_1\_INT} 的屏蔽中断状态。(RO)
\label{fielddesc:PMU0P2ACNTTARGET0REACH0HPINTST}\item [PMU\_0P2A\_CNT\_TARGET0\_REACH\_0\_HP\_INT\_ST] \hyperref[int:PMU0P2ACNTTARGET0REACH0INT]{PMU\_0P2A\_CNT\_TARGET0\_REACH\_0\_INT} 的屏蔽中断状态。(RO)
\label{fielddesc:PMU0P2ACNTTARGET1REACH0HPINTST}\item [PMU\_0P2A\_CNT\_TARGET1\_REACH\_0\_HP\_INT\_ST] \hyperref[int:PMU0P2ACNTTARGET1REACH0INT]{PMU\_0P2A\_CNT\_TARGET1\_REACH\_0\_INT} 的屏蔽中断状态。(RO)
\label{fielddesc:PMU0P2ACNTTARGET0REACH1HPINTST}\item [PMU\_0P2A\_CNT\_TARGET0\_REACH\_1\_HP\_INT\_ST] \hyperref[int:PMU0P2ACNTTARGET0REACH1INT]{PMU\_0P2A\_CNT\_TARGET0\_REACH\_1\_INT} 的屏蔽中断状态。(RO)
\label{fielddesc:PMU0P2ACNTTARGET1REACH1HPINTST}\item [PMU\_0P2A\_CNT\_TARGET1\_REACH\_1\_HP\_INT\_ST] \hyperref[int:PMU0P2ACNTTARGET1REACH1INT]{PMU\_0P2A\_CNT\_TARGET1\_REACH\_1\_INT} 的屏蔽中断状态。(RO)
\label{fielddesc:PMULPCPUEXCINTST}\item [PMU\_LP\_CPU\_EXC\_INT\_ST] \hyperref[int:PMULPCPUEXCINT]{PMU\_LP\_CPU\_EXC\_INT} 的屏蔽中断状态。(RO)
\label{fielddesc:PMUSDIOIDLEINTST}\item [PMU\_SDIO\_IDLE\_INT\_ST] \hyperref[int:PMUSDIOIDLEINT]{PMU\_SDIO\_IDLE\_INT} 的屏蔽中断状态。 (RO)
\label{fielddesc:PMUSWINTST}\item [PMU\_SW\_INT\_ST] \hyperref[int:PMUSWINT]{PMU\_SW\_INT} 的屏蔽中断状态。 (RO)
\label{fielddesc:PMUSOCSLEEPREJECTINTST}\item [PMU\_SOC\_SLEEP\_REJECT\_INT\_ST] \hyperref[int:PMUSOCSLEEPREJECTINT]{PMU\_SOC\_SLEEP\_REJECT\_INT} 的屏蔽中断状态。 (RO)
\label{fielddesc:PMUSOCWAKEUPINTST}\item [PMU\_SOC\_WAKEUP\_INT\_ST] \hyperref[int:PMUSOCWAKEUPINT]{PMU\_SOC\_WAKEUP\_INT} 的屏蔽中断状态。 (RO)
\end{reglist}\end{regdesc}
\end{register}


\begin{register}{H}{PMU\_HP\_INT\_ENA\_REG}{0x016C}\label{regdesc:PMUHPINTENAREG}
\regfield{PMU\_SOC\_WAKEUP\_INT\_ENA}{1}{31}{{0}}%
\regfield{PMU\_SOC\_SLEEP\_REJECT\_INT\_ENA}{1}{30}{{0}}%
\regfield{PMU\_SW\_INT\_ENA}{1}{29}{{0}}%
\regfield{PMU\_SDIO\_IDLE\_INT\_ENA}{1}{28}{{0}}%
\regfield{PMU\_LP\_CPU\_EXC\_INT\_ENA}{1}{27}{{0}}%
\regfield{(reserved)}{1}{26}{0}%
\regfield{reserved}{1}{25}{{0}}%
\regfield{reserved}{1}{24}{{0}}%
\regfield{reserved}{1}{23}{{0}}%
\regfield{reserved}{1}{22}{{0}}%
\regfield{PMU\_0P2A\_CNT\_TARGET1\_REACH\_1\_HP\_INT\_ENA}{1}{21}{{0}}%
\regfield{PMU\_0P2A\_CNT\_TARGET0\_REACH\_1\_HP\_INT\_ENA}{1}{20}{{0}}%
\regfield{PMU\_0P2A\_CNT\_TARGET1\_REACH\_0\_HP\_INT\_ENA}{1}{19}{{0}}%
\regfield{PMU\_0P2A\_CNT\_TARGET0\_REACH\_0\_HP\_INT\_ENA}{1}{18}{{0}}%
\regfield{PMU\_0P1A\_CNT\_TARGET1\_REACH\_1\_HP\_INT\_ENA}{1}{17}{{0}}%
\regfield{PMU\_0P1A\_CNT\_TARGET0\_REACH\_1\_HP\_INT\_ENA}{1}{16}{{0}}%
\regfield{PMU\_0P1A\_CNT\_TARGET1\_REACH\_0\_HP\_INT\_ENA}{1}{15}{{0}}%
\regfield{PMU\_0P1A\_CNT\_TARGET0\_REACH\_0\_HP\_INT\_ENA}{1}{14}{{0}}%
\regfield{(reserved)}{14}{0}{00000000000000}%
\reglabel{Reset}\regnewline%
\begin{regdesc}\begin{reglist}
\label{fielddesc:PMU0P1ACNTTARGET0REACH0HPINTENA}\item  [PMU\_0P1A\_CNT\_TARGET0\_REACH\_0\_HP\_INT\_ENA] 写 1 使能 \hyperref[int:PMU0P1ACNTTARGET0REACH0INT]{PMU\_0P1A\_CNT\_TARGET0\_REACH\_0\_INT}。(R/W)
\label{fielddesc:PMU0P1ACNTTARGET1REACH0HPINTENA}\item [PMU\_0P1A\_CNT\_TARGET1\_REACH\_0\_HP\_INT\_ENA] 写 1 使能 \hyperref[int:PMU0P1ACNTTARGET1REACH0INT]{PMU\_0P1A\_CNT\_TARGET1\_REACH\_0\_INT}。(R/W)
\label{fielddesc:PMU0P1ACNTTARGET0REACH1HPINTENA}\item [PMU\_0P1A\_CNT\_TARGET0\_REACH\_1\_HP\_INT\_ENA] 写 1 使能 \hyperref[int:PMU0P1ACNTTARGET0REACH1INT]{PMU\_0P1A\_CNT\_TARGET0\_REACH\_1\_INT}。(R/W)
\label{fielddesc:PMU0P1ACNTTARGET1REACH1HPINTENA}\item [PMU\_0P1A\_CNT\_TARGET1\_REACH\_1\_HP\_INT\_ENA] 写 1 使能 \hyperref[int:PMU0P1ACNTTARGET1REACH1INT]{PMU\_0P1A\_CNT\_TARGET1\_REACH\_1\_INT}。(R/W)
\label{fielddesc:PMU0P2ACNTTARGET0REACH0HPINTENA}\item [PMU\_0P2A\_CNT\_TARGET0\_REACH\_0\_HP\_INT\_ENA] 写 1 使能 \hyperref[int:PMU0P2ACNTTARGET0REACH0INT]{PMU\_0P2A\_CNT\_TARGET0\_REACH\_0\_INT}。(R/W)
\label{fielddesc:PMU0P2ACNTTARGET1REACH0HPINTENA}\item [PMU\_0P2A\_CNT\_TARGET1\_REACH\_0\_HP\_INT\_ENA] 写 1 使能 \hyperref[int:PMU0P2ACNTTARGET1REACH0INT]{PMU\_0P2A\_CNT\_TARGET1\_REACH\_0\_INT}。(R/W)
\label{fielddesc:PMU0P2ACNTTARGET0REACH1HPINTENA}\item [PMU\_0P2A\_CNT\_TARGET0\_REACH\_1\_HP\_INT\_ENA] 写 1 使能 \hyperref[int:PMU0P2ACNTTARGET0REACH1INT]{PMU\_0P2A\_CNT\_TARGET0\_REACH\_1\_INT}。(R/W)
\label{fielddesc:PMU0P2ACNTTARGET1REACH1HPINTENA}\item [PMU\_0P2A\_CNT\_TARGET1\_REACH\_1\_HP\_INT\_ENA] 写 1 使能 \hyperref[int:PMU0P2ACNTTARGET1REACH1INT]{PMU\_0P2A\_CNT\_TARGET1\_REACH\_1\_INT}。(R/W)
\label{fielddesc:PMULPCPUEXCINTENA}\item [PMU\_LP\_CPU\_EXC\_INT\_ENA] 写 1 使能 \hyperref[int:PMULPCPUEXCINT]{PMU\_LP\_CPU\_EXC\_INT}。(R/W)
\label{fielddesc:PMUSDIOIDLEINTENA}\item [PMU\_SDIO\_IDLE\_INT\_ENA] 写 1 使能 \hyperref[int:PMUSDIOIDLEINT]{PMU\_SDIO\_IDLE\_INT}。(R/W)
\label{fielddesc:PMUSWINTENA}\item [PMU\_SW\_INT\_ENA] 写 1 使能 \hyperref[int:PMUSWINT]{PMU\_SW\_INT}。(R/W)
\label{fielddesc:PMUSOCSLEEPREJECTINTENA}\item [PMU\_SOC\_SLEEP\_REJECT\_INT\_ENA] 写 1 使能 \hyperref[int:PMUSOCSLEEPREJECTINT]{PMU\_SOC\_SLEEP\_REJECT\_INT}。(R/W)
\label{fielddesc:PMUSOCWAKEUPINTENA}\item [PMU\_SOC\_WAKEUP\_INT\_ENA] 写 1 使能 \hyperref[int:PMUSOCWAKEUPINT]{PMU\_SOC\_WAKEUP\_INT}。  (R/W)
\end{reglist}\end{regdesc}
\end{register}


\begin{register}{H}{PMU\_HP\_INT\_CLR\_REG}{0x0170}\label{regdesc:PMUHPINTCLRREG}
\regfield{PMU\_SOC\_WAKEUP\_INT\_CLR}{1}{31}{{0}}%
\regfield{PMU\_SOC\_SLEEP\_REJECT\_INT\_CLR}{1}{30}{{0}}%
\regfield{PMU\_SW\_INT\_CLR}{1}{29}{{0}}%
\regfield{PMU\_SDIO\_IDLE\_INT\_CLR}{1}{28}{{0}}%
\regfield{PMU\_LP\_CPU\_EXC\_INT\_CLR}{1}{27}{{0}}%
\regfield{(reserved)}{5}{22}{00000}%
\regfield{PMU\_0P2A\_CNT\_TARGET1\_REACH\_1\_HP\_INT\_CLR}{1}{21}{{0}}%
\regfield{PMU\_0P2A\_CNT\_TARGET0\_REACH\_1\_HP\_INT\_CLR}{1}{20}{{0}}%
\regfield{PMU\_0P2A\_CNT\_TARGET1\_REACH\_0\_HP\_INT\_CLR}{1}{19}{{0}}%
\regfield{PMU\_0P2A\_CNT\_TARGET0\_REACH\_0\_HP\_INT\_CLR}{1}{18}{{0}}%
\regfield{PMU\_0P1A\_CNT\_TARGET1\_REACH\_1\_HP\_INT\_CLR}{1}{17}{{0}}%
\regfield{PMU\_0P1A\_CNT\_TARGET0\_REACH\_1\_HP\_INT\_CLR}{1}{16}{{0}}%
\regfield{PMU\_0P1A\_CNT\_TARGET1\_REACH\_0\_HP\_INT\_CLR}{1}{15}{{0}}%
\regfield{PMU\_0P1A\_CNT\_TARGET0\_REACH\_0\_HP\_INT\_CLR}{1}{14}{{0}}%
\regfield{(reserved)}{14}{0}{00000000000000}%
\reglabel{Reset}\regnewline%
\begin{regdesc}\begin{reglist}
\label{fielddesc:PMU0P1ACNTTARGET0REACH0HPINTCLR}\item  [PMU\_0P1A\_CNT\_TARGET0\_REACH\_0\_HP\_INT\_CLR] 写 1 清除 \hyperref[int:PMU0P1ACNTTARGET0REACH0INT]{PMU\_0P1A\_CNT\_TARGET0\_REACH\_0\_INT}。(WT)
\label{fielddesc:PMU0P1ACNTTARGET1REACH0HPINTCLR}\item [PMU\_0P1A\_CNT\_TARGET1\_REACH\_0\_HP\_INT\_CLR] 写 1 清除 \hyperref[int:PMU0P1ACNTTARGET1REACH0INT]{PMU\_0P1A\_CNT\_TARGET1\_REACH\_0\_INT}。(WT)
\label{fielddesc:PMU0P1ACNTTARGET0REACH1HPINTCLR}\item [PMU\_0P1A\_CNT\_TARGET0\_REACH\_1\_HP\_INT\_CLR] 写 1 清除 \hyperref[int:PMU0P1ACNTTARGET0REACH1INT]{PMU\_0P1A\_CNT\_TARGET0\_REACH\_1\_INT}。(WT)
\label{fielddesc:PMU0P1ACNTTARGET1REACH1HPINTCLR}\item [PMU\_0P1A\_CNT\_TARGET1\_REACH\_1\_HP\_INT\_CLR] 写 1 清除 \hyperref[int:PMU0P1ACNTTARGET1REACH1INT]{PMU\_0P1A\_CNT\_TARGET1\_REACH\_1\_INT}。(WT)
\label{fielddesc:PMU0P2ACNTTARGET0REACH0HPINTCLR}\item [PMU\_0P2A\_CNT\_TARGET0\_REACH\_0\_HP\_INT\_CLR] 写 1 清除 V\hyperref[int:PMU0P2ACNTTARGET0REACH0INT]{PMU\_0P2A\_CNT\_TARGET0\_REACH\_0\_INT}。(WT)
\label{fielddesc:PMU0P2ACNTTARGET1REACH0HPINTCLR}\item [PMU\_0P2A\_CNT\_TARGET1\_REACH\_0\_HP\_INT\_CLR] 写 1 清除 \hyperref[int:PMU0P2ACNTTARGET1REACH0INT]{PMU\_0P2A\_CNT\_TARGET1\_REACH\_0\_INT}。(WT)
\label{fielddesc:PMU0P2ACNTTARGET0REACH1HPINTCLR}\item [PMU\_0P2A\_CNT\_TARGET0\_REACH\_1\_HP\_INT\_CLR] 写 1 清除 \hyperref[int:PMU0P2ACNTTARGET0REACH1INT]{PMU\_0P2A\_CNT\_TARGET0\_REACH\_1\_INT}。(WT)
\label{fielddesc:PMU0P2ACNTTARGET1REACH1HPINTCLR}\item [PMU\_0P2A\_CNT\_TARGET1\_REACH\_1\_HP\_INT\_CLR] 写 1 清除 \hyperref[int:PMU0P2ACNTTARGET1REACH1INT]{PMU\_0P2A\_CNT\_TARGET1\_REACH\_1\_INT}。(WT)
\label{fielddesc:PMULPCPUEXCINTCLR}\item [PMU\_LP\_CPU\_EXC\_INT\_CLR] 写 1 清除 \hyperref[int:PMULPCPUEXCINT]{PMU\_LP\_CPU\_EXC\_INT}。 (WT)
\label{fielddesc:PMUSDIOIDLEINTCLR}\item [PMU\_SDIO\_IDLE\_INT\_CLR] 写 1 清除 \hyperref[int:PMUSDIOIDLEINT]{PMU\_SDIO\_IDLE\_INT}。 (WT)
\label{fielddesc:PMUSWINTCLR}\item [PMU\_SW\_INT\_CLR] 写 1 清除 \hyperref[int:PMUSWINT]{PMU\_SW\_INT}。(WT)
\label{fielddesc:PMUSOCSLEEPREJECTINTCLR}\item [PMU\_SOC\_SLEEP\_REJECT\_INT\_CLR] 写 1 清除 \hyperref[int:PMUSOCSLEEPREJECTINT]{PMU\_SOC\_SLEEP\_REJECT\_INT}。 (WT)
\label{fielddesc:PMUSOCWAKEUPINTCLR}\item [PMU\_SOC\_WAKEUP\_INT\_CLR] 写 1 清除 \hyperref[int:PMUSOCWAKEUPINT]{PMU\_SOC\_WAKEUP\_INT}。(WT)
\end{reglist}\end{regdesc}
\end{register}


\begin{register}{H}{PMU\_LP\_INT\_RAW\_REG}{0x0174}\label{regdesc:PMULPINTRAWREG}
\regfield{PMU\_HP\_SW\_TRIGGER\_INT\_RAW}{1}{31}{{0}}%
\regfield{PMU\_ACTIVE\_SWITCH\_SLEEP\_START\_INT\_RAW}{1}{30}{{0}}%
\regfield{PMU\_SLEEP\_SWITCH\_ACTIVE\_START\_INT\_RAW}{1}{29}{{0}}%
\regfield{PMU\_ACTIVE\_SWITCH\_SLEEP\_END\_INT\_RAW}{1}{28}{{0}}%
\regfield{PMU\_SLEEP\_SWITCH\_ACTIVE\_END\_INT\_RAW}{1}{27}{{0}}%
\regfield{PMU\_LP\_CPU\_WAKEUP\_INT\_RAW}{1}{26}{{0}}%
\regfield{(reserved)}{4}{22}{0000}%
\regfield{PMU\_0P2A\_CNT\_TARGET1\_REACH\_1\_LP\_INT\_RAW}{1}{21}{{0}}%
\regfield{PMU\_0P2A\_CNT\_TARGET0\_REACH\_1\_LP\_INT\_RAW}{1}{20}{{0}}%
\regfield{PMU\_0P2A\_CNT\_TARGET1\_REACH\_0\_LP\_INT\_RAW}{1}{19}{{0}}%
\regfield{PMU\_0P2A\_CNT\_TARGET0\_REACH\_0\_LP\_INT\_RAW}{1}{18}{{0}}%
\regfield{PMU\_0P1A\_CNT\_TARGET1\_REACH\_1\_LP\_INT\_RAW}{1}{17}{{0}}%
\regfield{PMU\_0P1A\_CNT\_TARGET0\_REACH\_1\_LP\_INT\_RAW}{1}{16}{{0}}%
\regfield{PMU\_0P1A\_CNT\_TARGET1\_REACH\_0\_LP\_INT\_RAW}{1}{15}{{0}}%
\regfield{PMU\_0P1A\_CNT\_TARGET0\_REACH\_0\_LP\_INT\_RAW}{1}{14}{{0}}%
\regfield{PMU\_LP\_CPU\_SLEEP\_REJECT\_INT\_RAW}{1}{13}{{0}}%
\regfield{(reserved)}{13}{0}{0000000000000}%
\reglabel{Reset}\regnewline%
\begin{regdesc}\begin{reglist}
\label{fielddesc:PMULPCPUSLEEPREJECTINTRAW}\item [PMU\_LP\_CPU\_SLEEP\_REJECT\_INT\_RAW] \hyperref[int:PMULPCPUSLEEPREJECTINT]{PMU\_LP\_CPU\_SLEEP\_REJECT\_INT} 的原始中断位。(R/WTC/SS)
\label{fielddesc:PMU0P1ACNTTARGET0REACH0LPINTRAW}\item [PMU\_0P1A\_CNT\_TARGET0\_REACH\_0\_LP\_INT\_RAW] \hyperref[int:PMU0P1ACNTTARGET0REACH0LPINT]{PMU\_0P1A\_CNT\_TARGET0\_REACH\_0\_LP\_INT} 的原始中断位。(R/WTC/SS)
\label{fielddesc:PMU0P1ACNTTARGET1REACH0LPINTRAW}\item [PMU\_0P1A\_CNT\_TARGET1\_REACH\_0\_LP\_INT\_RAW] \hyperref[int:PMU0P1ACNTTARGET1REACH0LPINT]{PMU\_0P1A\_CNT\_TARGET1\_REACH\_0\_LP\_INT} 的原始中断位。(R/WTC/SS)
\label{fielddesc:PMU0P1ACNTTARGET0REACH1LPINTRAW}\item [PMU\_0P1A\_CNT\_TARGET0\_REACH\_1\_LP\_INT\_RAW] \hyperref[int:PMU0P1ACNTTARGET0REACH1LPINT]{PMU\_0P1A\_CNT\_TARGET0\_REACH\_1\_LP\_INT} 的原始中断位。(R/WTC/SS)
\label{fielddesc:PMU0P1ACNTTARGET1REACH1LPINTRAW}\item [PMU\_0P1A\_CNT\_TARGET1\_REACH\_1\_LP\_INT\_RAW] \hyperref[int:PMU0P1ACNTTARGET1REACH1LPINT]{PMU\_0P1A\_CNT\_TARGET1\_REACH\_1\_LP\_INT} 的原始中断位。(R/WTC/SS)
\label{fielddesc:PMU0P2ACNTTARGET0REACH0LPINTRAW}\item [PMU\_0P2A\_CNT\_TARGET0\_REACH\_0\_LP\_INT\_RAW] \hyperref[int:PMU0P2ACNTTARGET0REACH0LPINT]{PMU\_0P2A\_CNT\_TARGET0\_REACH\_0\_LP\_INT} 的原始中断位。(R/WTC/SS)
\label{fielddesc:PMU0P2ACNTTARGET1REACH0LPINTRAW}\item [PMU\_0P2A\_CNT\_TARGET1\_REACH\_0\_LP\_INT\_RAW] \hyperref[int:PMU0P2ACNTTARGET1REACH0LPINT]{PMU\_0P2A\_CNT\_TARGET1\_REACH\_0\_LP\_INT} 的原始中断位。(R/WTC/SS)
\label{fielddesc:PMU0P2ACNTTARGET0REACH1LPINTRAW}\item [PMU\_0P2A\_CNT\_TARGET0\_REACH\_1\_LP\_INT\_RAW] \hyperref[int:PMU0P2ACNTTARGET0REACH1LPINT]{PMU\_0P2A\_CNT\_TARGET0\_REACH\_1\_LP\_INT} 的原始中断位。(R/WTC/SS)
\label{fielddesc:PMU0P2ACNTTARGET1REACH1LPINTRAW}\item [PMU\_0P2A\_CNT\_TARGET1\_REACH\_1\_LP\_INT\_RAW] \hyperref[int:PMU0P2ACNTTARGET1REACH1LPINT]{PMU\_0P2A\_CNT\_TARGET1\_REACH\_1\_LP\_INT} 的原始中断位。(R/WTC/SS)
\label{fielddesc:PMULPCPUWAKEUPINTRAW}\item [PMU\_LP\_CPU\_WAKEUP\_INT\_RAW] \hyperref[int:PMULPCPUWAKEUPINT]{PMU\_LP\_CPU\_WAKEUP\_INT} 的原始中断状态。 (R/WTC/SS)
\label{fielddesc:PMUSLEEPSWITCHACTIVEENDINTRAW}\item [PMU\_SLEEP\_SWITCH\_ACTIVE\_END\_INT\_RAW] \hyperref[int:PMUSLEEPSWITCHACTIVEENDINT]{PMU\_SLEEP\_SWITCH\_ACTIVE\_END\_INT} 的原始中断状态。 (R/WTC/SS)
\label{fielddesc:PMUACTIVESWITCHSLEEPENDINTRAW}\item [PMU\_ACTIVE\_SWITCH\_SLEEP\_END\_INT\_RAW] \hyperref[int:PMUACTIVESWITCHSLEEPENDINT]{PMU\_ACTIVE\_SWITCH\_SLEEP\_END\_INT} 的原始中断状态。(R/WTC/SS)
\label{fielddesc:PMUSLEEPSWITCHACTIVESTARTINTRAW}\item [PMU\_SLEEP\_SWITCH\_ACTIVE\_START\_INT\_RAW] \hyperref[int:PMUSLEEPSWITCHACTIVESTARTINT]{PMU\_SLEEP\_SWITCH\_ACTIVE\_START\_INT} 的原始中断状态。(R/WTC/SS)

\item [见下页...]
\end{reglist}\end{regdesc}
\end{register}

\addtocounter{Regfloat}{-1}
\begin{register}{H}{PMU\_LP\_INT\_RAW\_REG}{0x0174}
\begin{regdesc}\begin{reglist}
\item [接上页...]

\label{fielddesc:PMUACTIVESWITCHSLEEPSTARTINTRAW}\item [PMU\_ACTIVE\_SWITCH\_SLEEP\_START\_INT\_RAW] \hyperref[int:PMUACTIVESWITCHSLEEPSTARTINT]{PMU\_ACTIVE\_SWITCH\_SLEEP\_START\_INT} 的原始中断状态。 (R/WTC/SS)
\label{fielddesc:PMUHPSWTRIGGERINTRAW}\item [PMU\_HP\_SW\_TRIGGER\_INT\_RAW] \hyperref[int:PMUHPSWTRIGGERINT]{PMU\_HP\_SW\_TRIGGER\_INT} 的原始中断状态。(R/WTC/SS)
\end{reglist}\end{regdesc}
\end{register}


\begin{register}{H}{PMU\_LP\_INT\_ST\_REG}{0x0178}\label{regdesc:PMULPINTSTREG}
\regfield{PMU\_HP\_SW\_TRIGGER\_INT\_ST}{1}{31}{{0}}%
\regfield{PMU\_ACTIVE\_SWITCH\_SLEEP\_START\_INT\_ST}{1}{30}{{0}}%
\regfield{PMU\_SLEEP\_SWITCH\_ACTIVE\_START\_INT\_ST}{1}{29}{{0}}%
\regfield{PMU\_ACTIVE\_SWITCH\_SLEEP\_END\_INT\_ST}{1}{28}{{0}}%
\regfield{PMU\_SLEEP\_SWITCH\_ACTIVE\_END\_INT\_ST}{1}{27}{{0}}%
\regfield{PMU\_LP\_CPU\_WAKEUP\_INT\_ST}{1}{26}{{0}}%
\regfield{(reserved)}{4}{22}{0000}%
\regfield{PMU\_0P2A\_CNT\_TARGET1\_REACH\_1\_LP\_INT\_ST}{1}{21}{{0}}%
\regfield{PMU\_0P2A\_CNT\_TARGET0\_REACH\_1\_LP\_INT\_ST}{1}{20}{{0}}%
\regfield{PMU\_0P2A\_CNT\_TARGET1\_REACH\_0\_LP\_INT\_ST}{1}{19}{{0}}%
\regfield{PMU\_0P2A\_CNT\_TARGET0\_REACH\_0\_LP\_INT\_ST}{1}{18}{{0}}%
\regfield{PMU\_0P1A\_CNT\_TARGET1\_REACH\_1\_LP\_INT\_ST}{1}{17}{{0}}%
\regfield{PMU\_0P1A\_CNT\_TARGET0\_REACH\_1\_LP\_INT\_ST}{1}{16}{{0}}%
\regfield{PMU\_0P1A\_CNT\_TARGET1\_REACH\_0\_LP\_INT\_ST}{1}{15}{{0}}%
\regfield{PMU\_0P1A\_CNT\_TARGET0\_REACH\_0\_LP\_INT\_ST}{1}{14}{{0}}%
\regfield{PMU\_LP\_CPU\_SLEEP\_REJECT\_INT\_ST}{1}{13}{{0}}%
\regfield{(reserved)}{13}{0}{0000000000000}%
\reglabel{Reset}\regnewline%
\begin{regdesc}\begin{reglist}
\label{fielddesc:PMULPCPUSLEEPREJECTINTST}\item [PMU\_LP\_CPU\_SLEEP\_REJECT\_INT\_ST] \hyperref[int:PMULPCPUSLEEPREJECTINT]{PMU\_LP\_CPU\_SLEEP\_REJECT\_INT} 的原始中断位。(RO)
\label{fielddesc:PMU0P1ACNTTARGET0REACH0LPINTST}\item  [PMU\_0P1A\_CNT\_TARGET0\_REACH\_0\_LP\_INT\_ST] \hyperref[int:PMU0P1ACNTTARGET0REACH0LPINT]{PMU\_0P1A\_CNT\_TARGET0\_REACH\_0\_LP\_INT} 的屏蔽中断位。(RO)
\label{fielddesc:PMU0P1ACNTTARGET1REACH0LPINTST}\item [PMU\_0P1A\_CNT\_TARGET1\_REACH\_0\_LP\_INT\_ST] \hyperref[int:PMU0P1ACNTTARGET1REACH0LPINT]{PMU\_0P1A\_CNT\_TARGET1\_REACH\_0\_LP\_INT} 的屏蔽中断位。(RO)
\label{fielddesc:PMU0P1ACNTTARGET0REACH1LPINTST}\item [PMU\_0P1A\_CNT\_TARGET0\_REACH\_1\_LP\_INT\_ST] \hyperref[int:PMU0P1ACNTTARGET0REACH1LPINT]{PMU\_0P1A\_CNT\_TARGET0\_REACH\_1\_LP\_INT} 的屏蔽中断位。(RO)
\label{fielddesc:PMU0P1ACNTTARGET1REACH1LPINTST}\item [PMU\_0P1A\_CNT\_TARGET1\_REACH\_1\_LP\_INT\_ST] \hyperref[int:PMU0P1ACNTTARGET1REACH1LPINT]{PMU\_0P1A\_CNT\_TARGET1\_REACH\_1\_LP\_INT} 的屏蔽中断位。(RO)
\label{fielddesc:PMU0P2ACNTTARGET0REACH0LPINTST}\item [PMU\_0P2A\_CNT\_TARGET0\_REACH\_0\_LP\_INT\_ST] \hyperref[int:PMU0P2ACNTTARGET0REACH0LPINT]{PMU\_0P2A\_CNT\_TARGET0\_REACH\_0\_LP\_INT} 的屏蔽中断位。(RO)
\label{fielddesc:PMU0P2ACNTTARGET1REACH0LPINTST}\item [PMU\_0P2A\_CNT\_TARGET1\_REACH\_0\_LP\_INT\_ST] \hyperref[int:PMU0P2ACNTTARGET1REACH0LPINT]{PMU\_0P2A\_CNT\_TARGET1\_REACH\_0\_LP\_INT} 的屏蔽中断位。(RO)
\label{fielddesc:PMU0P2ACNTTARGET0REACH1LPINTST}\item [PMU\_0P2A\_CNT\_TARGET0\_REACH\_1\_LP\_INT\_ST] \hyperref[int:PMU0P2ACNTTARGET0REACH1LPINT]{PMU\_0P2A\_CNT\_TARGET0\_REACH\_1\_LP\_INT} 的屏蔽中断位。(RO)
\label{fielddesc:PMU0P2ACNTTARGET1REACH1LPINTST}\item [PMU\_0P2A\_CNT\_TARGET1\_REACH\_1\_LP\_INT\_ST] \hyperref[int:PMU0P2ACNTTARGET1REACH1LPINT]{PMU\_0P2A\_CNT\_TARGET1\_REACH\_1\_LP\_INT} 的屏蔽中断位。(RO)
\label{fielddesc:PMULPCPUWAKEUPINTST}\item [PMU\_LP\_CPU\_WAKEUP\_INT\_ST] \hyperref[int:PMULPCPUWAKEUPINT]{PMU\_LP\_CPU\_WAKEUP\_INT} 的屏蔽中断状态。 (RO)
\label{fielddesc:PMUSLEEPSWITCHACTIVEENDINTST}\item [PMU\_SLEEP\_SWITCH\_ACTIVE\_END\_INT\_ST] \hyperref[int:PMUSLEEPSWITCHACTIVEENDINT]{PMU\_SLEEP\_SWITCH\_ACTIVE\_END\_INT} 的屏蔽中断状态。(RO)
\label{fielddesc:PMUACTIVESWITCHSLEEPENDINTST}\item [PMU\_ACTIVE\_SWITCH\_SLEEP\_END\_INT\_ST] \hyperref[int:PMUACTIVESWITCHSLEEPENDINT]{PMU\_ACTIVE\_SWITCH\_SLEEP\_END\_INT} 的屏蔽中断状态。(RO)
\label{fielddesc:PMUSLEEPSWITCHACTIVESTARTINTST}\item [PMU\_SLEEP\_SWITCH\_ACTIVE\_START\_INT\_ST] \hyperref[int:PMUSLEEPSWITCHACTIVESTARTINT]{PMU\_SLEEP\_SWITCH\_ACTIVE\_START\_INT} 的屏蔽中断状态。 (RO)

\item [见下页...]
\end{reglist}\end{regdesc}
\end{register}

\addtocounter{Regfloat}{-1}
\begin{register}{H}{PMU\_LP\_INT\_ST\_REG}{0x0178}
\begin{regdesc}\begin{reglist}
\item [接上页...]

\label{fielddesc:PMUACTIVESWITCHSLEEPSTARTINTST}\item [PMU\_ACTIVE\_SWITCH\_SLEEP\_START\_INT\_ST] \hyperref[int:PMUACTIVESWITCHSLEEPSTARTINT]{PMU\_ACTIVE\_SWITCH\_SLEEP\_START\_INT} 的屏蔽中断状态。 (RO)
\label{fielddesc:PMUHPSWTRIGGERINTST}\item [PMU\_HP\_SW\_TRIGGER\_INT\_ST] \hyperref[int:PMUHPSWTRIGGERINT]{PMU\_HP\_SW\_TRIGGER\_INT} 的屏蔽中断状态。(RO)
\end{reglist}\end{regdesc}
\end{register}


\begin{register}{H}{PMU\_LP\_INT\_ENA\_REG}{0x017C}\label{regdesc:PMULPINTENAREG}
\regfield{PMU\_HP\_SW\_TRIGGER\_INT\_ENA}{1}{31}{{0}}%
\regfield{PMU\_ACTIVE\_SWITCH\_SLEEP\_START\_INT\_ENA}{1}{30}{{0}}%
\regfield{PMU\_SLEEP\_SWITCH\_ACTIVE\_START\_INT\_ENA}{1}{29}{{0}}%
\regfield{PMU\_ACTIVE\_SWITCH\_SLEEP\_END\_INT\_ENA}{1}{28}{{0}}%
\regfield{PMU\_SLEEP\_SWITCH\_ACTIVE\_END\_INT\_ENA}{1}{27}{{0}}%
\regfield{PMU\_LP\_CPU\_WAKEUP\_INT\_ENA}{1}{26}{{0}}%
\regfield{(reserved)}{4}{22}{0000}%
\regfield{PMU\_0P2A\_CNT\_TARGET1\_REACH\_1\_LP\_INT\_ENA}{1}{21}{{0}}%
\regfield{PMU\_0P2A\_CNT\_TARGET0\_REACH\_1\_LP\_INT\_ENA}{1}{20}{{0}}%
\regfield{PMU\_0P2A\_CNT\_TARGET1\_REACH\_0\_LP\_INT\_ENA}{1}{19}{{0}}%
\regfield{PMU\_0P2A\_CNT\_TARGET0\_REACH\_0\_LP\_INT\_ENA}{1}{18}{{0}}%
\regfield{PMU\_0P1A\_CNT\_TARGET1\_REACH\_1\_LP\_INT\_ENA}{1}{17}{{0}}%
\regfield{PMU\_0P1A\_CNT\_TARGET0\_REACH\_1\_LP\_INT\_ENA}{1}{16}{{0}}%
\regfield{PMU\_0P1A\_CNT\_TARGET1\_REACH\_0\_LP\_INT\_ENA}{1}{15}{{0}}%
\regfield{PMU\_0P1A\_CNT\_TARGET0\_REACH\_0\_LP\_INT\_ENA}{1}{14}{{0}}%
\regfield{PMU\_LP\_CPU\_SLEEP\_REJECT\_INT\_ENA}{1}{13}{{0}}%
\regfield{(reserved)}{13}{0}{0000000000000}%
\reglabel{Reset}\regnewline%
\begin{regdesc}\begin{reglist}
\label{fielddesc:PMULPCPUSLEEPREJECTINTENA}\item [PMU\_LP\_CPU\_SLEEP\_REJECT\_INT\_ENA] 写 1 使能 \hyperref[int:PMULPCPUSLEEPREJECTINT]{PMU\_LP\_CPU\_SLEEP\_REJECT\_INT}。(R/W)
\label{fielddesc:PMU0P1ACNTTARGET0REACH0LPINTENA}\item [PMU\_0P1A\_CNT\_TARGET0\_REACH\_0\_LP\_INT\_ENA] 写 1 使能 \hyperref[int:PMU0P1ACNTTARGET0REACH0LPINT]{PMU\_0P1A\_CNT\_TARGET0\_REACH\_0\_LP\_INT}。(R/W)
\label{fielddesc:PMU0P1ACNTTARGET1REACH0LPINTENA}\item [PMU\_0P1A\_CNT\_TARGET1\_REACH\_0\_LP\_INT\_ENA] 写 1 使能 \hyperref[int:PMU0P1ACNTTARGET1REACH0LPINT]{PMU\_0P1A\_CNT\_TARGET1\_REACH\_0\_LP\_INT}。(R/W)
\label{fielddesc:PMU0P1ACNTTARGET0REACH1LPINTENA}\item [PMU\_0P1A\_CNT\_TARGET0\_REACH\_1\_LP\_INT\_ENA] 写 1 使能 \hyperref[int:PMU0P1ACNTTARGET0REACH1LPINT]{PMU\_0P1A\_CNT\_TARGET0\_REACH\_1\_LP\_INT}。(R/W)
\label{fielddesc:PMU0P1ACNTTARGET1REACH1LPINTENA}\item [PMU\_0P1A\_CNT\_TARGET1\_REACH\_1\_LP\_INT\_ENA] 写 1 使能 \hyperref[int:PMU0P1ACNTTARGET1REACH1LPINT]{PMU\_0P1A\_CNT\_TARGET1\_REACH\_1\_LP\_INT}。(R/W)
\label{fielddesc:PMU0P2ACNTTARGET0REACH0LPINTENA}\item [PMU\_0P2A\_CNT\_TARGET0\_REACH\_0\_LP\_INT\_ENA] 写 1 使能 \hyperref[int:PMU0P2ACNTTARGET0REACH0LPINT]{PMU\_0P2A\_CNT\_TARGET0\_REACH\_0\_LP\_INT}。(R/W)
\label{fielddesc:PMU0P2ACNTTARGET1REACH0LPINTENA}\item [PMU\_0P2A\_CNT\_TARGET1\_REACH\_0\_LP\_INT\_ENA] 写 1 使能 \hyperref[int:PMU0P2ACNTTARGET1REACH0LPINT]{PMU\_0P2A\_CNT\_TARGET1\_REACH\_0\_LP\_INT}。(R/W)
\label{fielddesc:PMU0P2ACNTTARGET0REACH1LPINTENA}\item [PMU\_0P2A\_CNT\_TARGET0\_REACH\_1\_LP\_INT\_ENA] 写 1 使能 \hyperref[int:PMU0P2ACNTTARGET0REACH1LPINT]{PMU\_0P2A\_CNT\_TARGET0\_REACH\_1\_LP\_INT}。(R/W)
\label{fielddesc:PMU0P2ACNTTARGET1REACH1LPINTENA}\item [PMU\_0P2A\_CNT\_TARGET1\_REACH\_1\_LP\_INT\_ENA] 写 1 使能 \hyperref[int:PMU0P2ACNTTARGET1REACH1LPINT]{PMU\_0P2A\_CNT\_TARGET1\_REACH\_1\_LP\_INT}。(R/W)
\label{fielddesc:PMULPCPUWAKEUPINTENA}\item [PMU\_LP\_CPU\_WAKEUP\_INT\_ENA] 写 1 使能 \hyperref[int:PMULPCPUWAKEUPINT]{PMU\_LP\_CPU\_WAKEUP\_INT}。(R/W)
\label{fielddesc:PMUSLEEPSWITCHACTIVEENDINTENA}\item [PMU\_SLEEP\_SWITCH\_ACTIVE\_END\_INT\_ENA] 写 1 使能 \hyperref[int:PMUSLEEPSWITCHACTIVEENDINT]{PMU\_SLEEP\_SWITCH\_ACTIVE\_END\_INT}。(R/W)
\label{fielddesc:PMUACTIVESWITCHSLEEPENDINTENA}\item [PMU\_ACTIVE\_SWITCH\_SLEEP\_END\_INT\_ENA] 写 1 使能 \hyperref[int:PMUACTIVESWITCHSLEEPENDINT]{PMU\_ACTIVE\_SWITCH\_SLEEP\_END\_INT}。(R/W)
\label{fielddesc:PMUSLEEPSWITCHACTIVESTARTINTENA}\item [PMU\_SLEEP\_SWITCH\_ACTIVE\_START\_INT\_ENA] 写 1 使能 \hyperref[int:PMUSLEEPSWITCHACTIVESTARTINT]{PMU\_SLEEP\_SWITCH\_ACTIVE\_START\_INT}。(R/W)

\item [见下页...]
\end{reglist}\end{regdesc}
\end{register}

\addtocounter{Regfloat}{-1}
\begin{register}{H}{PMU\_LP\_INT\_ENA\_REG}{0x017C}
\begin{regdesc}\begin{reglist}
\item [接上页...]

\label{fielddesc:PMUACTIVESWITCHSLEEPSTARTINTENA}\item [PMU\_ACTIVE\_SWITCH\_SLEEP\_START\_INT\_ENA] 写 1 使能 \hyperref[int:PMUACTIVESWITCHSLEEPSTARTINT]{PMU\_ACTIVE\_SWITCH\_SLEEP\_START\_INT}。(R/W)
\label{fielddesc:PMUHPSWTRIGGERINTENA}\item [PMU\_HP\_SW\_TRIGGER\_INT\_ENA] 写 1 使能 \hyperref[int:PMUHPSWTRIGGERINT]{PMU\_HP\_SW\_TRIGGER\_INT}。(R/W)
\end{reglist}\end{regdesc}
\end{register}


\begin{register}{H}{PMU\_LP\_INT\_CLR\_REG}{0x0180}\label{regdesc:PMULPINTCLRREG}
\regfield{PMU\_HP\_SW\_TRIGGER\_INT\_CLR}{1}{31}{{0}}%
\regfield{PMU\_ACTIVE\_SWITCH\_SLEEP\_START\_INT\_CLR}{1}{30}{{0}}%
\regfield{PMU\_SLEEP\_SWITCH\_ACTIVE\_START\_INT\_CLR}{1}{29}{{0}}%
\regfield{PMU\_ACTIVE\_SWITCH\_SLEEP\_END\_INT\_CLR}{1}{28}{{0}}%
\regfield{PMU\_SLEEP\_SWITCH\_ACTIVE\_END\_INT\_CLR}{1}{27}{{0}}%
\regfield{PMU\_LP\_CPU\_WAKEUP\_INT\_CLR}{1}{26}{{0}}%
\regfield{(reserved)}{4}{22}{0000}%
\regfield{PMU\_0P2A\_CNT\_TARGET1\_REACH\_1\_LP\_INT\_CLR}{1}{21}{{0}}%
\regfield{PMU\_0P2A\_CNT\_TARGET0\_REACH\_1\_LP\_INT\_CLR}{1}{20}{{0}}%
\regfield{PMU\_0P2A\_CNT\_TARGET1\_REACH\_0\_LP\_INT\_CLR}{1}{19}{{0}}%
\regfield{PMU\_0P2A\_CNT\_TARGET0\_REACH\_0\_LP\_INT\_CLR}{1}{18}{{0}}%
\regfield{PMU\_0P1A\_CNT\_TARGET1\_REACH\_1\_LP\_INT\_CLR}{1}{17}{{0}}%
\regfield{PMU\_0P1A\_CNT\_TARGET0\_REACH\_1\_LP\_INT\_CLR}{1}{16}{{0}}%
\regfield{PMU\_0P1A\_CNT\_TARGET1\_REACH\_0\_LP\_INT\_CLR}{1}{15}{{0}}%
\regfield{PMU\_0P1A\_CNT\_TARGET0\_REACH\_0\_LP\_INT\_CLR}{1}{14}{{0}}%
\regfield{PMU\_LP\_CPU\_SLEEP\_REJECT\_LP\_INT\_CLR}{1}{13}{{0}}%
\regfield{(reserved)}{13}{0}{0000000000000}%
\reglabel{Reset}\regnewline%
\begin{regdesc}\begin{reglist}
\label{fielddesc:PMULPCPUSLEEPREJECTLPINTCLR}\item [PMU\_LP\_CPU\_SLEEP\_REJECT\_LP\_INT\_CLR] 写 1 清除 \hyperref[int:PMULPCPUSLEEPREJECTINT]{PMU\_LP\_CPU\_SLEEP\_REJECT\_INT}。(WT)
\label{fielddesc:PMU0P1ACNTTARGET0REACH0LPINTCLR}\item [PMU\_0P1A\_CNT\_TARGET0\_REACH\_0\_LP\_INT\_CLR] 写 1 清除 \hyperref[int:PMU0P1ACNTTARGET0REACH0LPINT]{PMU\_0P1A\_CNT\_TARGET0\_REACH\_0\_LP\_INT}。(WT)
\label{fielddesc:PMU0P1ACNTTARGET1REACH0LPINTCLR}\item [PMU\_0P1A\_CNT\_TARGET1\_REACH\_0\_LP\_INT\_CLR] 写 1 清除 \hyperref[int:PMU0P1ACNTTARGET1REACH0LPINT]{PMU\_0P1A\_CNT\_TARGET1\_REACH\_0\_LP\_INT}。(WT)
\label{fielddesc:PMU0P1ACNTTARGET0REACH1LPINTCLR}\item [PMU\_0P1A\_CNT\_TARGET0\_REACH\_1\_LP\_INT\_CLR] 写 1 清除 \hyperref[int:PMU0P1ACNTTARGET0REACH1LPINT]{PMU\_0P1A\_CNT\_TARGET0\_REACH\_1\_LP\_INT}。(WT)
\label{fielddesc:PMU0P1ACNTTARGET1REACH1LPINTCLR}\item [PMU\_0P1A\_CNT\_TARGET1\_REACH\_1\_LP\_INT\_CLR] 写 1 清除 \hyperref[int:PMU0P1ACNTTARGET1REACH1LPINT]{PMU\_0P1A\_CNT\_TARGET1\_REACH\_1\_LP\_INT}。(WT)
\label{fielddesc:PMU0P2ACNTTARGET0REACH0LPINTCLR}\item [PMU\_0P2A\_CNT\_TARGET0\_REACH\_0\_LP\_INT\_CLR] 写 1 清除 \hyperref[int:PMU0P2ACNTTARGET0REACH0LPINT]{PMU\_0P2A\_CNT\_TARGET0\_REACH\_0\_LP\_INT}。(WT)
\label{fielddesc:PMU0P2ACNTTARGET1REACH0LPINTCLR}\item [PMU\_0P2A\_CNT\_TARGET1\_REACH\_0\_LP\_INT\_CLR] 写 1 清除 \hyperref[int:PMU0P2ACNTTARGET1REACH0LPINT]{PMU\_0P2A\_CNT\_TARGET1\_REACH\_0\_LP\_INT}。(WT)
\label{fielddesc:PMU0P2ACNTTARGET0REACH1LPINTCLR}\item [PMU\_0P2A\_CNT\_TARGET0\_REACH\_1\_LP\_INT\_CLR] 写 1 清除 \hyperref[int:PMU0P2ACNTTARGET0REACH1LPINT]{PMU\_0P2A\_CNT\_TARGET0\_REACH\_1\_LP\_INT}。(WT)
\label{fielddesc:PMU0P2ACNTTARGET1REACH1LPINTCLR}\item [PMU\_0P2A\_CNT\_TARGET1\_REACH\_1\_LP\_INT\_CLR] 写 1 清除 \hyperref[int:PMU0P2ACNTTARGET1REACH1LPINT]{PMU\_0P2A\_CNT\_TARGET1\_REACH\_1\_LP\_INT}。(WT)
\label{fielddesc:PMULPCPUWAKEUPINTCLR}\item [PMU\_LP\_CPU\_WAKEUP\_INT\_CLR] 写 1 清除 \hyperref[int:PMULPCPUWAKEUPINT]{PMU\_LP\_CPU\_WAKEUP\_INT}。(WT)
\label{fielddesc:PMUSLEEPSWITCHACTIVEENDINTCLR}\item [PMU\_SLEEP\_SWITCH\_ACTIVE\_END\_INT\_CLR] 写 1 清除 \hyperref[int:PMUSLEEPSWITCHACTIVEENDINT]{PMU\_SLEEP\_SWITCH\_ACTIVE\_END\_INT}。(WT)
\label{fielddesc:PMUACTIVESWITCHSLEEPENDINTCLR}\item [PMU\_ACTIVE\_SWITCH\_SLEEP\_END\_INT\_CLR] 写 1 清除 \hyperref[int:PMUACTIVESWITCHSLEEPENDINT]{PMU\_ACTIVE\_SWITCH\_SLEEP\_END\_INT}。(WT)
\label{fielddesc:PMUSLEEPSWITCHACTIVESTARTINTCLR}\item [PMU\_SLEEP\_SWITCH\_ACTIVE\_START\_INT\_CLR] 写 1 清除 \hyperref[int:PMUSLEEPSWITCHACTIVESTARTINT]{PMU\_SLEEP\_SWITCH\_ACTIVE\_START\_INT}。(WT)

\item [见下页...]
\end{reglist}\end{regdesc}
\end{register}

\addtocounter{Regfloat}{-1}
\begin{register}{H}{PMU\_LP\_INT\_CLR\_REG}{0x0180}
\begin{regdesc}\begin{reglist}
\item [接上页...]

\label{fielddesc:PMUACTIVESWITCHSLEEPSTARTINTCLR}\item [PMU\_ACTIVE\_SWITCH\_SLEEP\_START\_INT\_CLR] 写 1 清除 \hyperref[int:PMUACTIVESWITCHSLEEPSTARTINT]{PMU\_ACTIVE\_SWITCH\_SLEEP\_START\_INT}。(WT)
\label{fielddesc:PMUHPSWTRIGGERINTCLR}\item [PMU\_HP\_SW\_TRIGGER\_INT\_CLR] 写 1 清除 \hyperref[int:PMUHPSWTRIGGERINT]{PMU\_HP\_SW\_TRIGGER\_INT}。 (WT)
\end{reglist}\end{regdesc}
\end{register}


\begin{register}{H}{PMU\_LP\_CPU\_PWR0\_REG}{0x0184}\label{regdesc:PMULPCPUPWR0REG}
\regfield{PMU\_LP\_CPU\_SLP\_BYPASS\_INTR\_EN}{1}{31}{{0}}%
\regfield{PMU\_LP\_CPU\_SLP\_RESET\_EN}{1}{30}{{0}}%
\regfield{PMU\_LP\_CPU\_SLP\_STALL\_EN}{1}{29}{{0}}%
\regfield{PMU\_LP\_CPU\_SLP\_STALL\_WAIT}{8}{21}{{0xff}}%
\regfield{PMU\_LP\_CPU\_SLP\_STALL\_FLAG\_EN}{1}{20}{{1}}%
\regfield{PMU\_LP\_CPU\_SLP\_WAITI\_FLAG\_EN}{1}{19}{{0}}%
\regfield{PMU\_LP\_CPU\_FORCE\_STALL}{1}{18}{{0}}%
\regfield{(reserved)}{16}{2}{0000000000000000}%
\regfield{PMU\_LP\_CPU\_STALL\_RDY}{1}{1}{{0}}%
\regfield{PMU\_LP\_CPU\_WAITI\_RDY}{1}{0}{{0}}%
\reglabel{Reset}\regnewline%
\begin{regdesc}\begin{reglist}
\label{fielddesc:PMULPCPUWAITIRDY}\item [PMU\_LP\_CPU\_WAITI\_RDY] 表示 LP CPU 是否进入 waiti 状态。\\
0：不进入状态\\
1：进入状态\\
(RO)
\label{fielddesc:PMULPCPUSTALLRDY}\item [PMU\_LP\_CPU\_STALL\_RDY] 表示 LP CPU 是否进入 stall 状态。\\
0：不进入状态\\
1：进入状态\\
(RO)
\label{fielddesc:PMULPCPUFORCESTALL}\item [PMU\_LP\_CPU\_FORCE\_STALL] 配置 LP CPU 进入 stall 状态。\\
0：不进入状态\\
1：进入状态\\
(R/W)
\label{fielddesc:PMULPCPUSLPWAITIFLAGEN}\item [PMU\_LP\_CPU\_SLP\_WAITI\_FLAG\_EN] 配置 LP CPU 在进入睡眠状态时是否需要进入 waiti 状态。\\
0：不进入状态\\
1：进入状态\\
 (R/W)\label{fielddesc:PMULPCPUSLPSTALLFLAGEN}\item [PMU\_LP\_CPU\_SLP\_STALL\_FLAG\_EN] 配置 LP CPU 在进入睡眠状态时是否需要进入 stall 状态。\\
0：不进入状态\\
1：进入状态\\
(R/W)\label{fielddesc:PMULPCPUSLPSTALLWAIT}\item [PMU\_LP\_CPU\_SLP\_STALL\_WAIT] 配置 LP CPU 在进入睡眠状态使能 stall 状态之后，等待 stall 生效的时间。单位是 LP\_DYN\_FAST\_CLK。(R/W)
\label{fielddesc:PMULPCPUSLPSTALLEN}\item [PMU\_LP\_CPU\_SLP\_STALL\_EN] 配置是否暂停进入睡眠模式的 LP CPU。\\
0：不暂停\\
1：暂停\\
(R/W)
\label{fielddesc:PMULPCPUSLPRESETEN}\item [PMU\_LP\_CPU\_SLP\_RESET\_EN] 配置是否复位进入睡眠模式的 LP CPU。\\
0：不复位\\
1：复位\\
(R/W)


\item [见下页...]
\end{reglist}\end{regdesc}
\end{register}

\addtocounter{Regfloat}{-1}
\begin{register}{H}{PMU\_LP\_CPU\_PWR0\_REG}{0x0184}
\begin{regdesc}\begin{reglist}
\item [接上页...]

\label{fielddesc:PMULPCPUSLPBYPASSINTREN}\item [PMU\_LP\_CPU\_SLP\_BYPASS\_INTR\_EN] 配置是否为进入睡眠模式的 LP CPU 打开中断信号。\\
0：关闭中断信号\\
1：打开中断信号\\
(R/W)
\end{reglist}\end{regdesc}
\end{register}


\begin{register}{H}{PMU\_LP\_CPU\_PWR1\_REG}{0x0188}\label{regdesc:PMULPCPUPWR1REG}
\regfield{PMU\_LP\_CPU\_SLEEP\_REQ}{1}{31}{{0}}%
\regfield{(reserved)}{31}{0}{0000000000000000000000000000000}%
\reglabel{Reset}\regnewline%
\begin{regdesc}\begin{reglist}
\label{fielddesc:PMULPCPUSLEEPREQ}\item [PMU\_LP\_CPU\_SLEEP\_REQ] LP CPU 睡眠请求寄存器。\\
0：不睡眠\\
1：进入睡眠\\
(WT)
\end{reglist}\end{regdesc}
\end{register}


\begin{register}{H}{PMU\_LP\_CPU\_PWR2\_REG}{0x018C}\label{regdesc:PMULPCPUPWR2REG}
\regfield{(reserved)}{1}{31}{0}%
\regfield{PMU\_LP\_CPU\_WAKEUP\_EN}{31}{0}{{0}}%
\reglabel{Reset}\regnewline%
\begin{regdesc}\begin{reglist}
\label{fielddesc:PMULPCPUWAKEUPEN}\item [PMU\_LP\_CPU\_WAKEUP\_EN] 配置 LP CPU 的唤醒源。具体请参考表 \ref{tab:lowpowm-wakeup-source} \textit{\nameref{tab:lowpowm-wakeup-source}}。(R/W)
\end{reglist}\end{regdesc}
\end{register}


\begin{register}{H}{PMU\_LP\_CPU\_PWR3\_REG}{0x0190}\label{regdesc:PMULPCPUPWR3REG}
\regfield{(reserved)}{1}{31}{0}%
\regfield{PMU\_LP\_CPU\_WAKEUP\_CAUSE}{31}{0}{{0}}%
\reglabel{Reset}\regnewline%
\begin{regdesc}\begin{reglist}
\label{fielddesc:PMULPCPUWAKEUPCAUSE}\item [PMU\_LP\_CPU\_WAKEUP\_CAUSE] 表示 LP CPU 唤醒的原因。(RO)
\end{reglist}\end{regdesc}
\end{register}


\begin{register}{H}{PMU\_LP\_CPU\_PWR4\_REG}{0x0194}\label{regdesc:PMULPCPUPWR4REG}
\regfield{(reserved)}{1}{31}{0}%
\regfield{PMU\_LP\_CPU\_REJECT\_EN}{31}{0}{{0}}%
\reglabel{Reset}\regnewline%
\begin{regdesc}\begin{reglist}
\label{fielddesc:PMULPCPUREJECTEN}\item [PMU\_LP\_CPU\_REJECT\_EN] 配置是否使能 LP CPU 拒绝睡眠功能。\\
0：禁能\\
1：使能\\
(R/W)
\end{reglist}\end{regdesc}
\end{register}


\begin{register}{H}{PMU\_LP\_CPU\_PWR5\_REG}{0x0198}\label{regdesc:PMULPCPUPWR5REG}
\regfield{(reserved)}{1}{31}{0}%
\regfield{PMU\_LP\_CPU\_REJECT\_CAUSE}{31}{0}{{0}}%
\reglabel{Reset}\regnewline%
\begin{regdesc}\begin{reglist}
\label{fielddesc:PMULPCPUREJECTCAUSE}\item [PMU\_LP\_CPU\_REJECT\_CAUSE] 表示 LP CPU 拒绝睡眠的原因。(RO)
\end{reglist}\end{regdesc}
\end{register}


\begin{register}{H}{PMU\_HP\_LP\_CPU\_COMM\_REG}{0x019C}\label{regdesc:PMUHPLPCPUCOMMREG}
\regfield{PMU\_HP\_TRIGGER\_LP}{1}{31}{{0}}%
\regfield{PMU\_LP\_TRIGGER\_HP}{1}{30}{{0}}%
\regfield{(reserved)}{30}{0}{000000000000000000000000000000}%
\reglabel{Reset}\regnewline%
\begin{regdesc}\begin{reglist}
\label{fielddesc:PMULPTRIGGERHP}\item [PMU\_LP\_TRIGGER\_HP] LP CPU 配置该寄存器为 1 时，芯片被唤醒。(WT)
\label{fielddesc:PMUHPTRIGGERLP}\item [PMU\_HP\_TRIGGER\_LP] HP CPU 配置该寄存器为 1 时，LP CPU 被唤醒。(WT)
\end{reglist}\end{regdesc}
\end{register}


\begin{register}{H}{PMU\_EXT\_LDO\_P0\_0P1A\_REG}{0x01B8}\label{regdesc:PMUEXTLDOP00P1AREG}
\regfield{reserved}{1}{31}{{0}}%
\regfield{PMU\_0P1A\_TARGET0\_0}{8}{23}{{0x80}}%
\regfield{PMU\_0P1A\_TARGET1\_0}{8}{15}{{0x40}}%
\regfield{PMU\_0P1A\_TIEH\_0}{1}{14}{{0}}%
\regfield{reserved}{1}{13}{{0}}%
\regfield{reserved}{1}{12}{{0}}%
\regfield{PMU\_0P1A\_TIEH\_SEL\_0}{3}{9}{{0}}%
\regfield{reserved}{1}{8}{{1}}%
\regfield{PMU\_0P1A\_FORCE\_TIEH\_SEL\_0}{1}{7}{{0}}%
\regfield{(reserved)}{7}{0}{0000000}%
\reglabel{Reset}\regnewline%
\begin{regdesc}\begin{reglist}
\label{fielddesc:PMU0P1AFORCETIEHSEL0}\item [PMU\_0P1A\_FORCE\_TIEH\_SEL\_0] 选择 VO1 调压器的模式选择控制权。 \\
1：模式选择由寄存器 \hyperref[fielddesc:PMU0P1ATIEHSEL0]{PMU\_0P1A\_TIEH\_SEL\_0} 控制 \\
0：模式选择由 eFuse 控制\\
(R/W)
\label{fielddesc:PMU0P1ATIEHSEL0}\item [PMU\_0P1A\_TIEH\_SEL\_0] 选择 VO1 调压器模式。\\
0：由寄存器 \hyperref[fielddesc:PMU0P1ATIEH0]{PMU\_0P1A\_TIEH\_0} 控制\\
1：由 SDMMC\_0 控制\\
2：透传模式\\
3：由 SDMMC\_1 控制\\
(R/W)
\label{fielddesc:PMU0P1ATIEH0}\item [PMU\_0P1A\_TIEH\_0] 选择 VO1 调压器模式。\\
0：LDO 模式\\
1：透传模式\\
(R/W)

\label{fielddesc:PMU0P1ATARGET10}\item [PMU\_0P1A\_TARGET1\_0] 配置 VO1 调压器等待计数器阶段 2 的超时时间，单位是 LP\_DYN\_SLOW\_CLK 的 16 分频。(R/W)
\label{fielddesc:PMU0P1ATARGET00}\item [PMU\_0P1A\_TARGET0\_0] 配置 VO1 调压器等待计数器阶段 1 的超时时间，单位是 LP\_DYN\_SLOW\_CLK 的 16 分频。(R/W)

\end{reglist}\end{regdesc}
\end{register}


\begin{register}{H}{PMU\_EXT\_LDO\_P0\_0P1A\_ANA\_REG}{0x01BC}\label{regdesc:PMUEXTLDOP00P1AANAREG}
\regfield{PMU\_ANA\_0P1A\_DREF\_0}{4}{28}{{0xb}}%
\regfield{PMU\_ANA\_0P1A\_EN\_CUR\_LIM\_0}{1}{27}{{0}}%
\regfield{PMU\_ANA\_0P1A\_EN\_VDET\_0}{1}{26}{{0}}%
\regfield{PMU\_ANA\_0P1A\_MUL\_0}{3}{23}{{0x2}}%
\regfield{(reserved)}{23}{0}{00000000000000000000000}%
\reglabel{Reset}\regnewline%
\begin{regdesc}\begin{reglist}
\label{fielddesc:PMUANA0P1AMUL0}\item [PMU\_ANA\_0P1A\_MUL\_0] 配置 MUL，调压器的输出电压为 MUL × Vref。\\
0-7：MUL 从 1 到 2.75，步进为 0.25。\\
(R/W)
\label{fielddesc:PMUANA0P1AENVDET0}\item [PMU\_ANA\_0P1A\_EN\_VDET\_0] 配置是否使能 VO1 电压检测。 \\
0：不使能 \\
1：使能 \\
(R/W)
\label{fielddesc:PMUANA0P1AENCURLIM0}\item [PMU\_ANA\_0P1A\_EN\_CUR\_LIM\_0] 配置是否使能 VO1 限流保护。\\
0：不使能 \\
1：使能 \\
(R/W)
\label{fielddesc:PMUANA0P1ADREF0}\item [PMU\_ANA\_0P1A\_DREF\_0] 配置 LDO Vref，调压器的输出电压为 MUL × Vref。\\
0-8：Vref = 0.5 V - 0.9 V，步进为 50 mV \\
9-15：Vref = 1.0 V - 1.6 V，步进为 100 mV\\
(R/W)
\end{reglist}\end{regdesc}
\end{register}


\begin{register}{H}{PMU\_EXT\_LDO\_P0\_0P2A\_REG}{0x01C0}\label{regdesc:PMUEXTLDOP00P2AREG}
\regfield{reserved}{1}{31}{{0}}%
\regfield{PMU\_0P2A\_TARGET0\_0}{8}{23}{{0x80}}%
\regfield{PMU\_0P2A\_TARGET1\_0}{8}{15}{{0x40}}%
\regfield{PMU\_0P2A\_TIEH\_0}{1}{14}{{0}}%
\regfield{reserved}{1}{13}{{0}}%
\regfield{reserved}{1}{12}{{0}}%
\regfield{PMU\_0P2A\_TIEH\_SEL\_0}{3}{9}{{0}}%
\regfield{PMU\_0P2A\_XPD\_0}{1}{8}{{1}}%
\regfield{PMU\_0P2A\_FORCE\_TIEH\_SEL\_0}{1}{7}{{0}}%
\regfield{(reserved)}{7}{0}{0000000}%
\reglabel{Reset}\regnewline%
\begin{regdesc}\begin{reglist}
\label{fielddesc:PMU0P2AFORCETIEHSEL0}\item [PMU\_0P2A\_FORCE\_TIEH\_SEL\_0] 选择 VO3 调压器模式。\\
0：由寄存器 \hyperref[fielddesc:PMU0P2ATIEHSEL0]{PMU\_0P2A\_TIEH\_SEL\_0} 控制\\
1：由 SDMMC\_0 控制\\
2：透传模式\\
3：由 SDMMC\_1 控制\\
(R/W)
\label{fielddesc:PMU0P2AXPD0}\item [PMU\_0P2A\_XPD\_0] 配置软件开启 VO3 调压器。\\
0：关闭 \\
1：打开 \\
(R/W)
\label{fielddesc:PMU0P2ATIEHSEL0}\item [PMU\_0P2A\_TIEH\_SEL\_0] 选择 VO3 调压器模式。\\
0：由寄存器 \hyperref[fielddesc:PMU0P2ATIEH0]{PMU\_0P2A\_TIEH\_0} 控制\\
1：由 SDMMC\_0 控制\\
2：透传模式\\
3：由 SDMMC\_1 控制\\
(R/W)
\label{fielddesc:PMU0P2ATIEH0}\item [PMU\_0P2A\_TIEH\_0] 选择 VO3 调压器模式。\\
0：LDO 模式\\
1：透传模式\\
(R/W)
\label{fielddesc:PMU0P2ATARGET10}\item [PMU\_0P2A\_TARGET1\_0] 配置 VO3 调压器等待计数器阶段 2 的超时时间，单位是 LP\_DYN\_SLOW\_CLK 的 16 分频。(R/W)
\label{fielddesc:PMU0P2ATARGET00}\item [PMU\_0P2A\_TARGET0\_0] 配置 VO3 调压器等待计数器阶段 1 的超时时间，单位是 LP\_DYN\_SLOW\_CLK 的 16 分频。(R/W)
\end{reglist}\end{regdesc}
\end{register}


\begin{register}{H}{PMU\_EXT\_LDO\_P0\_0P2A\_ANA\_REG}{0x01C4}\label{regdesc:PMUEXTLDOP00P2AANAREG}
\regfield{PMU\_ANA\_0P2A\_DREF\_0}{4}{28}{{0xb}}%
\regfield{PMU\_ANA\_0P2A\_EN\_CUR\_LIM\_0}{1}{27}{{0}}%
\regfield{PMU\_ANA\_0P2A\_EN\_VDET\_0}{1}{26}{{0}}%
\regfield{PMU\_ANA\_0P2A\_MUL\_0}{3}{23}{{0x2}}%
\regfield{(reserved)}{23}{0}{00000000000000000000000}%
\reglabel{Reset}\regnewline%
\begin{regdesc}\begin{reglist}
\label{fielddesc:PMUANA0P2AMUL0}\item [PMU\_ANA\_0P2A\_MUL\_0] 配置 MUL，调压器的输出电压为 MUL × Vref。\\
0-7：MUL 从 1 到 2.75，步进为 0.25。\\
(R/W)
\label{fielddesc:PMUANA0P2AENVDET0}\item [PMU\_ANA\_0P2A\_EN\_VDET\_0] 配置是否使能 VO3 电压检测。 \\
0：不使能 \\
1：使能 \\
(R/W)
\label{fielddesc:PMUANA0P2AENCURLIM0}\item [PMU\_ANA\_0P2A\_EN\_CUR\_LIM\_0] 配置是否使能 VO3 限流保护。\\
0：不使能 \\
1：使能 \\
(R/W)
\label{fielddesc:PMUANA0P2ADREF0}\item [PMU\_ANA\_0P2A\_DREF\_0] 配置 LDO Vref，调压器的输出电压为 MUL × Vref。\\
0-8：Vref = 0.5 V - 0.9 V，步进为 50 mV \\
9-15：Vref = 1.0 V - 1.6 V，步进为 100 mV\\
(R/W)
\end{reglist}\end{regdesc}
\end{register}








\begin{register}{H}{PMU\_EXT\_LDO\_P1\_0P1A\_REG}{0x01D0}\label{regdesc:PMUEXTLDOP10P1AREG}
\regfield{reserved}{1}{31}{{0}}%
\regfield{PMU\_0P1A\_TARGET0\_1}{8}{23}{{0x80}}%
\regfield{PMU\_0P1A\_TARGET1\_1}{8}{15}{{0x40}}%
\regfield{PMU\_0P1A\_TIEH\_1}{1}{14}{{0}}%
\regfield{reserved}{1}{13}{{0}}%
\regfield{reserved}{1}{12}{{0}}%
\regfield{PMU\_0P1A\_TIEH\_SEL\_1}{3}{9}{{0}}%
\regfield{PMU\_0P1A\_XPD\_1}{1}{8}{{1}}%
\regfield{PMU\_0P1A\_FORCE\_TIEH\_SEL\_1}{1}{7}{{0}}%
\regfield{(reserved)}{7}{0}{0000000}%
\reglabel{Reset}\regnewline%
\begin{regdesc}\begin{reglist}
\label{fielddesc:PMU0P1AFORCETIEHSEL1}\item [PMU\_0P1A\_FORCE\_TIEH\_SEL\_1] 选择 VO2 调压器模式。\\
0：由寄存器 \hyperref[fielddesc:PMU0P1ATIEHSEL1]{PMU\_0P1A\_TIEH\_SEL\_1} 控制\\
1：由 SDMMC\_0 控制\\
2：透传模式\\
3：由 SDMMC\_1 控制\\
(R/W)
\label{fielddesc:PMU0P1AXPD0}\item [PMU\_0P1A\_XPD\_1] 配置软件开启 VO2 调压器。\\
0：关闭 \\
1：打开 \\
(R/W)
\label{fielddesc:PMU0P1ATIEHSEL1}\item [PMU\_0P1A\_TIEH\_SEL\_1] 选择 VO2 调压器模式。\\
0：由寄存器 \hyperref[fielddesc:PMU0P1ATIEH1]{PMU\_0P1A\_TIEH\_1} 控制\\
1：由 SDMMC\_0 控制\\
2：透传模式\\
3：由 SDMMC\_1 控制\\
(R/W)
\label{fielddesc:PMU0P1ATIEH1}\item [PMU\_0P1A\_TIEH\_1] 选择 VO2 调压器模式。\\
0：LDO 模式\\
1：透传模式\\
(R/W)
\label{fielddesc:PMU0P1ATARGET11}\item [PMU\_0P1A\_TARGET1\_1] 配置 VO2 调压器等待计数器阶段 2 的超时时间，单位是 LP\_DYN\_SLOW\_CLK 的 16 分频。(R/W)\label{fielddesc:PMU0P1ATARGET01}\item [PMU\_0P1A\_TARGET0\_1] 配置 VO2 调压器等待计数器阶段 1 的超时时间，单位是 LP\_DYN\_SLOW\_CLK 的 16 分频。(R/W)
\end{reglist}\end{regdesc}
\end{register}


\begin{register}{H}{PMU\_EXT\_LDO\_P1\_0P1A\_ANA\_REG}{0x01D4}\label{regdesc:PMUEXTLDOP10P1AANAREG}
\regfield{PMU\_ANA\_0P1A\_DREF\_1}{4}{28}{{0xb}}%
\regfield{PMU\_ANA\_0P1A\_EN\_CUR\_LIM\_1}{1}{27}{{0}}%
\regfield{PMU\_ANA\_0P1A\_EN\_VDET\_1}{1}{26}{{0}}%
\regfield{PMU\_ANA\_0P1A\_MUL\_1}{3}{23}{{0x2}}%
\regfield{(reserved)}{23}{0}{00000000000000000000000}%
\reglabel{Reset}\regnewline%
\begin{regdesc}\begin{reglist}
\label{fielddesc:PMUANA0P1AMUL1}\item [PMU\_ANA\_0P1A\_MUL\_1] 配置 MUL，调压器的输出电压为 MUL × Vref。\\
0-7：MUL 从 1 到 2.75，步进为 0.25。\\
(R/W)
\label{fielddesc:PMUANA0P1AENVDET1}\item [PMU\_ANA\_0P1A\_EN\_VDET\_1] 配置是否使能 VO2 电压检测。\\
0：不使能 \\
1：使能 \\
(R/W)
\label{fielddesc:PMUANA0P1AENCURLIM1}\item [PMU\_ANA\_0P1A\_EN\_CUR\_LIM\_1] 配置是否使能 VO2 限流保护。\\
0：不使能 \\
1：使能 \\
(R/W)
\label{fielddesc:PMUANA0P1ADREF1}\item [PMU\_ANA\_0P1A\_DREF\_1] 配置 LDO Vref，调压器的输出电压为 MUL × Vref。\\
0-8：Vref = 0.5 V - 0.9 V，步进为 50 mV \\
9-15：Vref = 1.0 V - 1.6 V，步进为 100 mV\\
(R/W)

\end{reglist}\end{regdesc}
\end{register}


\begin{register}{H}{PMU\_EXT\_LDO\_P1\_0P2A\_REG}{0x01D8}\label{regdesc:PMUEXTLDOP10P2AREG}
\regfield{reserved}{1}{31}{{0}}%
\regfield{PMU\_0P2A\_TARGET0\_1}{8}{23}{{0x80}}%
\regfield{PMU\_0P2A\_TARGET1\_1}{8}{15}{{0x40}}%
\regfield{PMU\_0P2A\_TIEH\_1}{1}{14}{{0}}%
\regfield{reserved}{1}{13}{{0}}%
\regfield{reserved}{1}{12}{{0}}%
\regfield{PMU\_0P2A\_TIEH\_SEL\_1}{3}{9}{{0}}%
\regfield{PMU\_0P2A\_XPD\_1}{1}{8}{{1}}%
\regfield{PMU\_0P2A\_FORCE\_TIEH\_SEL\_1}{1}{7}{{0}}%
\regfield{(reserved)}{7}{0}{0000000}%
\reglabel{Reset}\regnewline%
\begin{regdesc}\begin{reglist}
\label{fielddesc:PMU0P2AFORCETIEHSEL1}\item [PMU\_0P2A\_FORCE\_TIEH\_SEL\_1] 选择 VO4 调压器模式。\\
0：由寄存器 \hyperref[fielddesc:PMU0P2ATIEHSEL1]{PMU\_0P2A\_TIEH\_SEL\_1} 控制\\
1：由 SDMMC\_0 控制\\
2：透传模式\\
3：由 SDMMC\_1 控制\\
(R/W)
 \label{fielddesc:PMU0P2AXPD1}\item [PMU\_0P2A\_XPD\_1] 配置软件开启 VO4 调压器。\\
0：关闭 \\
1：打开 \\
(R/W)
\label{fielddesc:PMU0P2ATIEHSEL1}\item [PMU\_0P2A\_TIEH\_SEL\_1] 选择 VO4 调压器模式。\\
0：由寄存器 \hyperref[fielddesc:PMU0P2ATIEH1]{PMU\_0P2A\_TIEH\_1} 控制\\
1：由 SDMMC\_0 控制\\
2：透传模式\\
3：由 SDMMC\_1 控制\\
(R/W)
\label{fielddesc:PMU0P2ATIEH1}\item [PMU\_0P2A\_TIEH\_1] 选择 VO4 调压器模式。\\
0：LDO 模式\\
1：透传模式\\
(R/W)
\label{fielddesc:PMU0P2ATARGET11}\item [PMU\_0P2A\_TARGET1\_1] 配置 VO4 调压器等待计数器阶段 2 的超时时间，单位是 LP\_DYN\_SLOW\_CLK 的 16 分频。(R/W)
\label{fielddesc:PMU0P2ATARGET01}\item [PMU\_0P2A\_TARGET0\_1]配置 VO4 调压器等待计数器阶段 1 的超时时间，单位是 LP\_DYN\_SLOW\_CLK 的 16 分频。(R/W)
\end{reglist}\end{regdesc}
\end{register}


\begin{register}{H}{PMU\_EXT\_LDO\_P1\_0P2A\_ANA\_REG}{0x01DC}\label{regdesc:PMUEXTLDOP10P2AANAREG}
\regfield{PMU\_ANA\_0P2A\_DREF\_1}{4}{28}{{0xb}}%
\regfield{PMU\_ANA\_0P2A\_EN\_CUR\_LIM\_1}{1}{27}{{0}}%
\regfield{PMU\_ANA\_0P2A\_EN\_VDET\_1}{1}{26}{{0}}%
\regfield{PMU\_ANA\_0P2A\_MUL\_1}{3}{23}{{0x2}}%
\regfield{(reserved)}{23}{0}{00000000000000000000000}%
\reglabel{Reset}\regnewline%
\begin{regdesc}\begin{reglist}
\label{fielddesc:PMUANA0P2AMUL1}\item [PMU\_ANA\_0P2A\_MUL\_1] 配置 MUL，调压器的输出电压为 MUL × Vref。\\
0-7：MUL 从 1 到 2.75，步进为 0.25。\\
(R/W)
\label{fielddesc:PMUANA0P2AENVDET1}\item [PMU\_ANA\_0P2A\_EN\_VDET\_1] 配置是否使能 VO4 电压检测。 \\
0：不使能 \\
1：使能 \\
(R/W)
\label{fielddesc:PMUANA0P2AENCURLIM1}\item [PMU\_ANA\_0P2A\_EN\_CUR\_LIM\_1] 配置是否使能 VO4 限流保护。 \\
0：不使能 \\
1：使能 \\
(R/W)
\label{fielddesc:PMUANA0P2ADREF1}\item [PMU\_ANA\_0P2A\_DREF\_1] 配置 LDO Vref，调压器的输出电压为 MUL × Vref。\\
0-8：Vref = 0.5 V - 0.9 V，步进为 50 mV \\
9-15：Vref = 1.0 V - 1.6 V，步进为 100 mV\\
(R/W)
\end{reglist}\end{regdesc}
\end{register}


\begin{register}{H}{PMU\_EXT\_WAKEUP\_LV\_REG}{0x01E8}\label{regdesc:PMUEXTWAKEUPLVREG}
\regfield{PMU\_EXT\_WAKEUP\_LV}{32}{0}{{0}}%
\reglabel{Reset}\regnewline%
\begin{regdesc}\begin{reglist}
\label{fielddesc:PMUEXTWAKEUPLV}\item [PMU\_EXT\_WAKEUP\_LV] 配置 EXT 唤醒的唤醒电平位图。\\
0：低电平唤醒 \\
1：高电平唤醒 \\
(R/W)
\end{reglist}\end{regdesc}
\end{register}


\begin{register}{H}{PMU\_EXT\_WAKEUP\_SEL\_REG}{0x01EC}\label{regdesc:PMUEXTWAKEUPSELREG}
\regfield{PMU\_EXT\_WAKEUP\_SEL}{32}{0}{{0}}%
\reglabel{Reset}\regnewline%
\begin{regdesc}\begin{reglist}
\label{fielddesc:PMUEXTWAKEUPSEL}\item [PMU\_EXT\_WAKEUP\_SEL] 配置 EXT 唤醒的唤醒使能位图。\\
0：禁止 \\
1：使能 \\
(R/W)

\end{reglist}\end{regdesc}
\end{register}


\begin{register}{H}{PMU\_EXT\_WAKEUP\_ST\_REG}{0x01F0}\label{regdesc:PMUEXTWAKEUPSTREG}
\regfield{PMU\_EXT\_WAKEUP\_STATUS}{32}{0}{{0}}%
\reglabel{Reset}\regnewline%
\begin{regdesc}\begin{reglist}
\label{fielddesc:PMUEXTWAKEUPSTATUS}\item [PMU\_EXT\_WAKEUP\_STATUS] 获取唤醒的 GPIO 编号。0 对应 GPIO0，1 对应 GPIO1，以此类推。(RO)
\end{reglist}\end{regdesc}
\end{register}


\begin{register}{H}{PMU\_EXT\_WAKEUP\_CNTL\_REG}{0x01F4}\label{regdesc:PMUEXTWAKEUPCNTLREG}
\regfield{PMU\_EXT\_WAKEUP\_FILTER}{1}{31}{{0}}%
\regfield{PMU\_EXT\_WAKEUP\_STATUS\_CLR}{1}{30}{{0}}%
\regfield{(reserved)}{30}{0}{000000000000000000000000000000}%
\reglabel{Reset}\regnewline%
\begin{regdesc}\begin{reglist}
\label{fielddesc:PMUEXTWAKEUPSTATUSCLR}\item [PMU\_EXT\_WAKEUP\_STATUS\_CLR] 清空唤醒记录。(R/W)
\label{fielddesc:PMUEXTWAKEUPFILTER}\item [PMU\_EXT\_WAKEUP\_FILTER] 配置是否使能唤醒滤波器。\\
0：禁止 \\
1：使能 \\
(R/W)
\end{reglist}\end{regdesc}
\end{register}


\begin{register}{H}{PMU\_SDIO\_WAKEUP\_CNTL\_REG}{0x01F8}\label{regdesc:PMUSDIOWAKEUPCNTLREG}
\regfield{(reserved)}{22}{10}{0000000000000000000000}%
\regfield{PMU\_SDIO\_ACT\_DNUM}{10}{0}{{0x3ff}}%
\reglabel{Reset}\regnewline%
\begin{regdesc}\begin{reglist}
\label{fielddesc:PMUSDIOACTDNUM}\item [PMU\_SDIO\_ACT\_DNUM] 配置 SDIO 的工作状态的维持时间。单位是 LP\_DYN\_FAST\_CLK。(R/W)
\end{reglist}\end{regdesc}
\end{register}


\begin{register}{H}{PMU\_TOUCH\_PWR\_CNTL\_REG}{0x0210}\label{regdesc:PMUTOUCHPWRCNTLREG}
\regfield{PMU\_TOUCH\_SLEEP\_TIMER\_EN}{1}{31}{{0}}%
\regfield{PMU\_TOUCH\_FORCE\_DONE}{1}{30}{{0}}%
\regfield{PMU\_TOUCH\_SLEEP\_CYCLES}{16}{14}{{100}}%
\regfield{PMU\_TOUCH\_WAIT\_CYCLES}{9}{5}{{10}}%
\regfield{(reserved)}{5}{0}{00000}%
\reglabel{Reset}\regnewline%
\begin{regdesc}\begin{reglist}
\label{fielddesc:PMUTOUCHWAITCYCLES}\item [PMU\_TOUCH\_WAIT\_CYCLES] 配置 TOUCH 定时器唤醒 PMU 等待的时间，单位是 LP\_DYN\_SLOW\_CLK。(R/W)
\label{fielddesc:PMUTOUCHSLEEPCYCLES}\item [PMU\_TOUCH\_SLEEP\_CYCLES] 配置 TOUCH 的睡眠时间，单位是 LP\_DYN\_SLOW\_CLK。(R/W)
\label{fielddesc:PMUTOUCHFORCEDONE}\item [PMU\_TOUCH\_FORCE\_DONE] 强制返回 TOUCH 定时器的工作状态。(R/W)
\label{fielddesc:PMUTOUCHSLEEPTIMEREN}\item [PMU\_TOUCH\_SLEEP\_TIMER\_EN] 使能 TOUCH 定时器开始工作。(R/W)
\end{reglist}\end{regdesc}
\end{register}



\begin{register}{H}{PMU\_DATE\_REG}{0x03FC}\label{regdesc:PMUDATEREG}
\regfield{PMU\_CLK\_EN}{1}{31}{{0}}%
\regfield{PMU\_PMU\_DATE}{31}{0}{{0x2303140}}%
\reglabel{Reset}\regnewline%
\begin{regdesc}\begin{reglist}
\label{fielddesc:PMUPMUDATE}\item [PMU\_PMU\_DATE] 版本寄存器。(R/W)
\label{fielddesc:PMUCLKEN}\item [PMU\_CLK\_EN] 强制开启寄存器的自动时钟门控。(R/W)
\end{reglist}\end{regdesc}
\end{register}



